% Options for packages loaded elsewhere
\PassOptionsToPackage{unicode}{hyperref}
\PassOptionsToPackage{hyphens}{url}
\documentclass[
]{article}
\usepackage{xcolor}
\usepackage[margin=25mm]{geometry}
\usepackage{amsmath,amssymb}
\setcounter{secnumdepth}{-\maxdimen} % remove section numbering
\usepackage{iftex}
\ifPDFTeX
  \usepackage[T1]{fontenc}
  \usepackage[utf8]{inputenc}
  \usepackage{textcomp} % provide euro and other symbols
\else % if luatex or xetex
  \usepackage{unicode-math} % this also loads fontspec
  \defaultfontfeatures{Scale=MatchLowercase}
  \defaultfontfeatures[\rmfamily]{Ligatures=TeX,Scale=1}
\fi
\usepackage{lmodern}
\ifPDFTeX\else
  % xetex/luatex font selection
\fi
% Use upquote if available, for straight quotes in verbatim environments
\IfFileExists{upquote.sty}{\usepackage{upquote}}{}
\IfFileExists{microtype.sty}{% use microtype if available
  \usepackage[]{microtype}
  \UseMicrotypeSet[protrusion]{basicmath} % disable protrusion for tt fonts
}{}
\makeatletter
\@ifundefined{KOMAClassName}{% if non-KOMA class
  \IfFileExists{parskip.sty}{%
    \usepackage{parskip}
  }{% else
    \setlength{\parindent}{0pt}
    \setlength{\parskip}{6pt plus 2pt minus 1pt}}
}{% if KOMA class
  \KOMAoptions{parskip=half}}
\makeatother
\setlength{\emergencystretch}{3em} % prevent overfull lines
\providecommand{\tightlist}{%
  \setlength{\itemsep}{0pt}\setlength{\parskip}{0pt}}
\usepackage{mathrsfs}
\usepackage{bookmark}
\IfFileExists{xurl.sty}{\usepackage{xurl}}{} % add URL line breaks if available
\urlstyle{same}
\hypersetup{
  pdftitle={Volume-independent local estimates in the full SU(2) Wilson vacuum},
  hidelinks,
  pdfcreator={LaTeX via pandoc}}

\title{Volume-independent local estimates in the full SU(2) Wilson
vacuum}
\author{}
\date{8 September 2026}

\begin{document}
\maketitle

\section{1. Original Hamiltonian, physical state, and geometric
input}\label{original-hamiltonian-physical-state-and-geometric-input}

Fix an integer \(L\ge2\). The vertices are
\(\mathsf V_L=\{-L,\ldots,L\}^3\). A physical link \(e=(n,i)\) is
oriented from \(n\) to \(n+\mathbf e_i\) whenever both endpoints are
vertices. Its independent variable is \(U_e\in SU(2)\); an inverse
traversal contributes \(U_e^{-1}\). Every contained elementary face
\(p=(n;i,j)\), \(i<j\), is included, with \[
W_p=\operatorname{tr}\left[
 U_i(n)U_j(n+\mathbf e_i)U_i(n+\mathbf e_j)^{-1}U_j(n)^{-1}
\right].                                                    \tag{1.1}
\] Write \(\mathsf E_L,\mathsf P_L\) for these complete sets, and \[
N=3(2L)(2L+1)^2,\qquad M=3(2L)^2(2L+1).
\] The Hilbert space is \(\mathcal H_L=L^2(SU(2)^{\mathsf E_L},dU)\),
with product Haar probability measure and inner product conjugate-linear
in its first entry. The differential operators and all physical
constants are \[
T_a=-i\sigma_a/2,\quad
X_{e,a}f=\left.\frac{d}{dt}f(\ldots,e^{tT_a}U_e,\ldots)\right|_{t=0},
\quad E_e=-\sum_{a=1}^3X_{e,a}^2,
\] \[
H=\kappa\sum_{e\in\mathsf E_L}E_e
       +b\sum_{p\in\mathsf P_L}(2-W_p),\qquad
\kappa=\frac{2g_{\rm YM}^2}{a},\quad
b=\frac{1}{2g_{\rm YM}^2a},\quad
\xi=\frac b\kappa=\frac1{4g_{\rm YM}^4},\quad a,g_{\rm YM}>0.       \tag{1.2}
\] In the original metric \(c(X,Y)=-\operatorname{tr}(XY)/2\), the
orthonormal frame is \(2T_a\), so \(\Delta_e^c=4\Delta_e^T\). Thus the
electric term in (1.2) also equals
\(-g_{\rm YM}^2\sum_e\Delta_e^c/(2a)\). No coefficient or energy origin
is changed. In particular the scalar \(2bM\) remains in \(H\).

The original geometric condition remains \[
\operatorname{Im}\tau
\left(\operatorname{Im}\beta-
\frac{6(\operatorname{Im}\mu)^2}{\operatorname{Im}\tau}\right)\ne0.
                                                               \tag{1.3}
\] The input magnetic patch uses
\(L_T=\operatorname{Im}\tau_T,\ q_T=-\operatorname{Im}\beta_T,
m_T=\operatorname{Im}\mu_T,\ D=L_Tq_T+6m_T^2>0\) and
\(\theta=2\pi a^2/D\). Its exact projected small-angle state, proved in
the parent source, is \[
\chi_\theta=-\frac23\theta^2 v_{w^{\rm nat}}+O_{H^1}(\theta^4),
\qquad w^{\rm nat}_{(n,1)}=n_2^2,\quad
w^{\rm nat}_{(n,2)}=w^{\rm nat}_{(n,3)}=0.                         \tag{1.4}
\] The angle remainder in (1.4) is a fixed-regulator input, not a bound
claimed uniform here. The result below concerns its actual electric
state with the exact original weights.

The product Laplacian on this compact connected manifold has domain
\(H^2\), form domain \(H^1\), and compact resolvent. Its smooth bounded
potential preserves these domains. A lowest form minimizer can be chosen
nonnegative by the gradient inequality for the modulus, obtained at its
zeros by approximation with \(\sqrt{|f|^2+\epsilon^2}\). Elliptic
regularity and the strong maximum principle make it smooth and strictly
positive. Write it as \(\psi>0\), with \(\int\psi^2dU=1\), and write its
physical energy as \(\mathcal E\). Integration by parts using
\(H\psi=\mathcal E\psi\) gives \[
\mathfrak q_{H-\mathcal E}[\psi F]
 =\kappa\sum_{e,a}\int\psi^2|X_{e,a}F|^2dU.                    \tag{1.5}
\] A second ground eigenfunction divided by \(\psi\) would have all
derivatives zero in (1.5), hence be constant. This proves simplicity.
The gauge action \(U_e\mapsto g_{s(e)}U_eg_{t(e)}^{-1}\) preserves the
operator, form, and positive unit vector; consequently \(\psi\) is gauge
invariant. All states below are formed from this same vacuum.

For each link set \[
r_e=\#\{p:e\in\partial p\},\quad
W_e^{\star}=\sum_{p\ni e}W_p,\quad
t_{ef}=\#\{p:e,f\in\partial p\},\quad
\gamma_e=\langle\psi,E_e\psi\rangle,\quad v_e=(E_e-\gamma_e)\psi.
\] Here \(2\le r_e\le4\), including boundary links. Define the actual
matrices \[
C_{ef}=\langle v_e,v_f\rangle,\qquad
\mathcal N_{ef}=\langle v_e,(H-\mathcal E)v_f\rangle.             \tag{1.6}
\] They are real symmetric positive-semidefinite matrices. Smoothness of
\(\psi\) supplies every operator and form domain used in (1.6).

\section{2. A local projection and an exact comparison
operator}\label{a-local-projection-and-an-exact-comparison-operator}

Let \(S\subset\mathsf E_L\) be a nonempty set of links. It need not be
connected. Define \[
\mathsf J(S)=\{p:\partial p\cap S\ne\varnothing\},\quad
m_S=|\mathsf J(S)|,\quad
W_S=\sum_{p\in\mathsf J(S)}W_p,\quad H_S=\sum_{e\in S}E_e,
\] \[
H_{\rm ext,S}=\kappa\sum_{e\notin S}E_e
       +b\sum_{p\notin\mathsf J(S)}(2-W_p).                    \tag{2.1}
\] The last operator acts on the exterior tensor factor only. Its bottom
energy is \(\mathcal E_{\rm ext,S}\); if that factor is empty it is the
zero operator on \(\mathbb C\). The exact decomposition is \[
H=\kappa H_S+H_{\rm ext,S}+2bm_S-bW_S.                          \tag{2.2}
\] Every face meeting \(S\), including faces crossing its boundary, is
present in \(W_S\). Every other face is present in \(H_{\rm ext,S}\).
The constant terms total \(2bM\) in both (1.2) and (2.2).

The map \[
(P_S f)(U_{S^c})=\int f(U_S,U_{S^c})\,dU_S,\qquad Q_S=I-P_S    \tag{2.3}
\] is understood with its output lifted as a function constant on \(S\).
Fubini proves \(P_S=P_S^*=P_S^2\). Its range is
\(\mathbb C1_S\otimes L^2(SU(2)^{S^c})\), and its kernel consists of
functions of zero conditional Haar integral on \(S\). It commutes with
all link Casimirs and with exterior operators, and preserves smooth
functions and Sobolev domains. Each plaquette in \(W_S\) contains a link
of \(S\) exactly once in the fundamental or inverse fundamental
representation. Integration of that link is zero: the substitution
\(U\mapsto -U\) changes its matrix entries' signs. Therefore \[
P_SW_SP_S=0.                                                   \tag{2.4}
\] This is an operator identity, not a statement that the interacting
vacuum has a Haar marginal. Haar invariance under endpoint left/right
multiplication also proves that \(P_S,Q_S\) intertwine the gauge action.

Take as a trial vector the constant on \(S\) tensored with an exterior
ground vector. Equation (2.4) makes its \(W_S\) expectation zero. The
full variational principle gives \[
\mathcal E\le\mathcal E_{\rm ext,S}+2bm_S.
\] Consequently the exterior self-adjoint operator \[
K_S=H_{\rm ext,S}+2bm_S-\mathcal E\quad\hbox{satisfies }K_S\ge0.
                                                               \tag{2.5}
\] No gap of \(K_S\) is assumed or needed. Put \(B_S=\kappa H_S+K_S\).
It is a sum of nonnegative operators on different tensor factors. Their
spectral resolutions strongly commute; the operator domain of their sum
is their domain intersection. The full smooth vacuum belongs to that
domain, and its exact equation is \[
B_S\psi=bW_S\psi.                                             \tag{2.6}
\] This identity retains the original interaction and exterior vacuum
dependence. It is not a free-vacuum eigen-equation.

The single-link eigenvalues of \(E_e\) are \(j(j+1)\),
\(j=0,\frac12,1,\ldots\), with the constant as its unique zero vector.
Indeed the Pauli matrices give \(-\sum_aT_a^2=3I/4\) on the fundamental
representation, and the spin raising/lowering matrices give \(j(j+1)\)
on each higher representation. Completeness of their matrix coefficients
under product Haar measure gives the joint spectral decomposition. Thus
\(H_S\ge(3/4)Q_S\). On \(Q_S\mathcal H_L\), \(B_S\ge3\kappa/4\). Define
\[
R_S=Q_S(B_S|_{Q_S\mathcal H_L})^{-1}Q_S.
\] Joint spectral calculus gives, without a Neumann expansion, \[
\|R_S\|\le\frac4{3\kappa},\qquad
\|E_eR_S\|\le\frac1\kappa\quad(e\in S).                      \tag{2.7}
\] For the latter assertion the multiplier is
\(\lambda_e/(\kappa\sum_{f\in S}\lambda_f+k)\le1/\kappa\) for
\(\lambda_f,k\ge0\). This also proves that the displayed composition has
a bounded extension on the full Hilbert space. Applying \(Q_S\) to (2.6)
now gives the exact local-resolvent equation \[
Q_S\psi=bR_SW_S\psi.                                          \tag{2.8}
\]

\section{3. Local electric moments at every positive
coupling}\label{local-electric-moments-at-every-positive-coupling}

\textbf{Theorem 3.1.} For every box \(L\ge2\), every \(a,g_{\rm YM}>0\),
every link \(e\), and every nonempty finite link set \(S\), \[
\|E_e\psi\|\le2r_e\xi,\qquad
0\le\gamma_e\le\frac{16}{3}r_e^2\xi^2,                         \tag{3.1}
\] \[
\|Q_S\psi\|\le q_S\xi,\qquad
q_S=\frac83\left(\sum_{e\in S}r_e^2\right)^{1/2}
\le\frac{32}{3}\sqrt{|S|}.                                    \tag{3.2}
\] The constants have no dependence on the number of exterior links or
faces.

\textbf{Proof.} Use \(S=\{e\}\) in (2.6). Its local interaction
satisfies \(\|W_S\|\le2r_e\). For the joint nonnegative spectral
variables of \(\kappa E_e\) and \(K_S\), their sum is at least the first
variable. Squaring this scalar inequality and integrating against the
spectral measure of \(\psi\) proves \[
\kappa\|E_e\psi\|\le\|B_S\psi\|
 =b\|W_e^{\star}\psi\|\le2br_e.
\] Since each nonzero eigenvalue of \(E_e\) is at least \(3/4\), the
spectral inequality \(E_e\le(4/3)E_e^2\) gives the second bound in
(3.1). The commuting projections onto nonconstant individual link
factors obey \(Q_S\le\sum_{e\in S}Q_{\{e\}}\) in their joint zero/one
eigenvalues. Also \(Q_{\{e\}}\le(16/9)E_e^2\). Taking expectations and
using (3.1) gives \[
\|Q_S\psi\|^2\le\frac{16}{9}\sum_{e\in S}\|E_e\psi\|^2
 \le\frac{64}{9}\xi^2\sum_{e\in S}r_e^2.
\] This proves (3.2). All estimates used the actual normalized vacuum;
no estimate on its full-box distance from the constant was used.
\(\square\)

For later use, let \(F\) be any bounded multiplication operator
depending only on \(S\), with full Haar mean \(\overline F\). With
\(\eta=P_S\psi,\ u=Q_S\psi,\ q=\|u\|\), Fubini gives
\(\langle\eta,F\eta\rangle=\overline F\|\eta\|^2\). Expanding all four
terms of \(\langle\eta+u,F(\eta+u)\rangle\) and using
\(\|\eta\|^2=1-q^2\) therefore gives \[
\left|\langle\psi,F\psi\rangle-\overline F\right|
 \le2\|F\|q+2\|F\|q^2.                                      \tag{3.3}
\] This is a proved bound on a marginal expectation. It does not replace
the marginal by Haar measure.

\section{4. An edge-star resolvent with an explicit state-vector
remainder}\label{an-edge-star-resolvent-with-an-explicit-state-vector-remainder}

Take \[
S_e=\bigcup_{p\ni e}\partial p,\qquad m_e^{\star}=m_{S_e}.
\] Each face at \(e\) adds at most three other links, so
\(|S_e|\le1+3r_e\le13\). Each link meets at most four faces, hence
\(m_e^{\star}\le4|S_e|\le52\). Equation (3.2) gives
\(q_{S_e}\le32\sqrt{13}/3<39\). These counts hold also at every boundary
of the open box.

\textbf{Theorem 4.1.} Simultaneously for all links and all boxes, on the
explicit interval \(0<\xi\le1/64\), \[
\left\|E_e\psi-\frac\xi4W_e^{\star}\psi\right\|
 \le4200\xi^2.                                                 \tag{4.1}
\] This compares two vectors in the same full interacting Hilbert space.

\textbf{Proof.} In Section 2 take \(S=S_e\) and abbreviate its operators
as \(P,Q,K,B,R\). Write \(\eta=P\psi,\ u=Q\psi\), with
\(\|u\|\le q_S\xi\). Applying \(P\) to (2.6) and using (2.4) gives \[
K\eta=bPW_Su,\qquad \|K\eta\|\le2bm_Sq_S\xi.                 \tag{4.2}
\] Since \(E_eP=0\), equation (2.8) gives \[
E_e\psi=bE_eRW_S\eta+bE_eRW_Su.                               \tag{4.3}
\] Its second term has norm at most \(2m_Sq_S\xi^2\), by (2.7). For a
face not containing \(e\), \(E_eW_p\eta=0\). For a face containing
\(e\), all four boundary links lie in \(S\), and \[
E_eW_p\eta=\frac34W_p\eta,\qquad H_SW_p\eta=3W_p\eta.
\] Here the four fundamental Casimirs contribute separately, with
exterior variables untouched. Multiplication by such \(W_p\) commutes
with \(K\). The bounded map \(\zeta\mapsto W_p\zeta\) from the exterior
space to the local four-fundamental sector intertwines \(3\kappa+K\)
with \(B\). Its image is in \(Q\mathcal H_L\). Resolvent intertwining
and (4.3) thus give the exact identity \[
bE_eRW_S\eta=\frac{3b}{4}
 W_e^{\star}(3\kappa+K)^{-1}\eta.
\] The exterior inverse exists because \(K\ge0\). Subtracting
\((\xi/4)W_e^{\star}\eta\), while retaining \(K\), gives \[
-\frac\xi4W_e^{\star}(3\kappa+K)^{-1}K\eta,
\] whose norm by (4.2) is at most \(r_em_Sq_S\xi^3/3\). Replacing
\(\eta\) by \(\psi\) in the comparison term has the additional norm cost
\(r_eq_S\xi^2/2\). We have proved the more detailed bound \[
\left\|E_e\psi-\frac\xi4W_e^{\star}\psi\right\|
\le (2m_S+r_e/2)q_S\xi^2
       +\frac{r_em_Sq_S}{3}\xi^3
\le4134\xi^2+2704\xi^3.                                      \tag{4.4}
\] For \(\xi\le1/64\) the last coefficient is
\(4134+2704/64=16705/4<4200\). This proves (4.1). Equations (4.2)--(4.4)
exhibit every discarded-in-an-estimate term as an exact retained term
with a bound; none is set to zero. \(\square\)

\section{5. Uniform local Wilson and connected electric
covariances}\label{uniform-local-wilson-and-connected-electric-covariances}

For elementary faces \(p,q\), Haar integration gives
\(\int W_pW_qdU=\delta_{pq}\). For unequal faces a link occurs in only
one support; its central sign change annihilates the integral. For equal
faces Haar invariance reduces the integral to
\(\int_{SU(2)}(\operatorname{tr}U)^2dU=1\). In quaternion coordinates
the latter is \(4\int u_0^2=1\), since the four coordinate second
moments on the unit three-sphere equal \(1/4\).

The function \(W_pW_q\) has norm at most \(4\) and uses at most eight
links. Equations (3.2)--(3.3), with \(q_S\le32\sqrt8/3<31\), prove, for
\(0<\xi\le1/64\), \[
|\langle W_pW_q\rangle_\psi-\delta_{pq}|
 \le248\xi+7688\xi^2\le370\xi.                              \tag{5.1}
\] This remains valid for faces at any separation; no finite-range
assertion about their true correlations is imposed.

For an individual face choose any one boundary link \(e\) and decompose
\(\psi=P_{\{e\}}\psi+Q_{\{e\}}\psi\). The first-first matrix element of
\(W_p\) vanishes by integration over \(e\). Equation (3.2) gives
\(q\le32\xi/3\), and the other three terms give \[
|\langle W_p\rangle_\psi|\le4q+2q^2
 \le\left(\frac{128}{3}+\frac{32}{9}\right)\xi<47\xi.
\] Consequently the connected Wilson moment obeys \[
|\langle W_pW_q\rangle_\psi
 -\langle W_p\rangle_\psi\langle W_q\rangle_\psi-\delta_{pq}|
 \le410\xi.                                                   \tag{5.2}
\] Indeed (5.1) and \(47^2\xi^2\le(2209/64)\xi\) give a smaller constant
than \(410\).

\textbf{Theorem 5.1.} On the same interval, simultaneously for all boxes
and links, \[
\left|C_{ef}-\frac{\xi^2}{16}t_{ef}\right|\le42500\xi^3.       \tag{5.3}
\]

\textbf{Proof.} Set \(u_e=E_e\psi\) and \(a_e=(\xi/4)W_e^{\star}\psi\).
The preceding bounds give \[
\|u_e\|\le8\xi,\quad \|a_e\|\le2\xi,\quad
\|u_e-a_e\|\le4200\xi^2.
\] Using the exact difference
\(\langle u_e-a_e,u_f\rangle+\langle a_e,u_f-a_f\rangle\), its absolute
value is at most \(42000\xi^3\). Also \[
\langle a_e,a_f\rangle
 =\frac{\xi^2}{16}\sum_{p\ni e,q\ni f}\langle W_pW_q\rangle_\psi.
\] There are at most sixteen ordered pairs in this sum. Its difference
from \(\xi^2t_{ef}/16\) is at most \(370\xi^3\), by (5.1). Finally (3.1)
bounds the exact centering term by \[
|\gamma_e\gamma_f|\le\frac{65536}{9}\xi^4
 \le\frac{1024}{9}\xi^3<114\xi^3.
\] The sum \(42000+370+114=42484<42500\) proves (5.3). No odd Taylor
coefficient has been asserted nonzero: this is an absolute real-coupling
remainder bound. The fixed-box evenness theorem is compatible with
(5.3); without an additional uniform analytic argument it does not
improve this bound to order \(\xi^4\). \(\square\)

\section{6. Energy entries and a volume-independent weighted numerator
estimate}\label{energy-entries-and-a-volume-independent-weighted-numerator-estimate}

Put \(A=H-\mathcal E\). The product rule, with the original sign of
\(E_f\), gives on the smooth vacuum \[
AE_f\psi=[H,E_f]\psi
=b\left(\frac34W_f^{\star}\psi-2D_f\psi\right),\qquad
D_f\psi=\sum_{p\ni f,a}(X_{f,a}W_p)X_{f,a}\psi.                \tag{6.1}
\] The electric sum commutes with \(E_f\). The potential contribution is
\(b[E_f,W_f^{\star}]\), retaining its first-derivative term in full.

For any one occurrence of a link in a simple Wilson loop, cyclically
basing the holonomy at that occurrence writes its derivative as
\(\operatorname{tr}(T_aU)\), or its negative, with a possible adjoint
rotation of the index \(a\). Quaternion coordinates then give \[
\sum_a|X_{f,a}W_p|^2=1-W_p^2/4\le1.                           \tag{6.2}
\] Pointwise Cauchy--Schwarz, followed by integration, bounds each
plaquette contraction in \(D_f\psi\) by
\((\sum_a\|X_{f,a}\psi\|^2)^{1/2}=\sqrt{\gamma_f}\). Thus (3.1) gives \[
\|D_f\psi\|\le r_f\sqrt{\gamma_f}
 \le\frac4{\sqrt3}r_f^2\xi\le\frac{64}{\sqrt3}\xi<37\xi.     \tag{6.3}
\]

\textbf{Theorem 6.1.} For \(0<\xi\le1/64\), all boxes, and all links, \[
\left|\mathcal N_{ef}-\frac{3\kappa\xi^2}{16}t_{ef}\right|
 \le27000\kappa\xi^3.                                        \tag{6.4}
\] For every real edge weight \(w\), including weights depending on
\(L\) and \(\xi\), \[
\left|w^T\mathcal Nw-\frac{3\kappa\xi^2}{16}\|\mathsf Iw\|_2^2\right|
 \le351000\kappa\xi^3\|w\|_2^2,\qquad
 (\mathsf Iw)_p=\sum_{e\in\partial p}w_e.                     \tag{6.5}
\]

\textbf{Proof.} Since \(A\psi=0\), centering does not change
\(\mathcal N_{ef}=\langle E_e\psi,AE_f\psi\rangle\). Insert (6.1).
Replacing \(E_e\psi\) by \(a_e\) in the first term costs at most
\(6\cdot4200\kappa\xi^3=25200\kappa\xi^3\), because
\(\|W_f^{\star}\psi\|\le8\). The comparison of its remaining Wilson
products to Haar costs at most \(1110\kappa\xi^3\) by (5.1). The full
derivative term is at most \(2\kappa\xi(8\xi)(37\xi)=592\kappa\xi^3\).
Their sum is \(26902\kappa\xi^3<27000\kappa\xi^3\), proving (6.4).

To prove the weighted conclusion without extending the entry bound to an
unjustified operator bound, expand the double commutator: \[
\mathcal N_{ef}
 =\tfrac12\langle\psi,[E_e,[H,E_f]]\psi\rangle
 =\frac b2\sum_{p:e,f\in\partial p}
 \langle\psi,[E_e,[2-W_p,E_f]]\psi\rangle.                     \tag{6.6}
\] The first identity follows from commutation of \(E_e,E_f\), the
vacuum equation at both endpoints, and reality of the two equal cross
terms. In the second, a face absent from \(f\) has zero inner
commutator. A face containing \(f\) but absent from \(e\) has
coefficients and derivatives commuting with \(E_e\). This proves exact
support on shared faces. A link has at most \(1+3r_e\le13\) partners,
counting itself. The remainder in (6.4) has the same support. For its
real symmetric matrix \(R\), \[
|w^TRw|\le\sum_{e,f}|R_{ef}|\frac{w_e^2+w_f^2}{2}
 \le13\cdot27000\kappa\xi^3\sum_ew_e^2.
\] Finally \(t=\mathsf I^T\mathsf I\) by its definition. This proves
(6.5), with no factor of \(M\), \(N\), or \(L\) in its constant.
\(\square\)

\section{7. Actual states, arbitrary weights, and the entire leading
kernel}\label{actual-states-arbitrary-weights-and-the-entire-leading-kernel}

For real \(w\) write
\(\Gamma_w=\sum_ew_eE_e,\ v_w=(\Gamma_w-\langle\Gamma_w\rangle_\psi)\psi\).
These operators and states preserve the gauge constraint; all products
are taken on smooth \(\psi\). Equations (5.3) and (6.5) give
simultaneously \[
\left|\|v_w\|^2-\frac{\xi^2}{16}\|\mathsf Iw\|_2^2\right|
 \le42500\xi^3\|w\|_1^2,                                    \tag{7.1}
\] \[
\left|\mathfrak q_A[v_w]-\frac{3\kappa\xi^2}{16}\|\mathsf Iw\|_2^2\right|
 \le351000\kappa\xi^3\|w\|_2^2.                              \tag{7.2}
\] The \(\ell^1\) norm in (7.1) is not replaced by an \(\ell^2\) norm.
The long-distance covariance entries are retained. Equations
(7.1)--(7.2) are genuine full-vacuum inequalities at fixed nonzero
coupling, not coefficients interpreted outside a finite-box perturbation
disk. Section 12 proves an additional \(\ell^2\) upper bound by a
convergent full-vacuum cluster construction. Sections 13--14
subsequently prove absolute distance decay and the stronger two-sided
\(\ell^2\) remainder (14.5). Equation (7.1) remains valid on its
original larger interval.

On the entire real kernel \(\mathsf Iw=0\), a sharper order statement
follows from the vector calculation. Summing (4.1) gives exactly
\(\Gamma_w\psi=\sum_ew_e(E_e\psi-a_e)\), since
\(\sum_ew_eW_e^{\star}=\sum_p(\mathsf Iw)_pW_p=0\). Projection
perpendicular to \(\psi\) is a contraction, so \[
\|v_w\|^2\le17640000\xi^4\|w\|_1^2.                          \tag{7.3}
\] In the sum of (6.1), the same cancellation gives
\(A\Gamma_w\psi=-2b\sum_ew_eD_e\psi\). Using (6.3) and (4.1), without
dividing by a possibly small state norm, proves \[
0\le\mathfrak q_A[v_w]\le310800\kappa\xi^4\|w\|_1^2.         \tag{7.4}
\] The constants are \(4200^2=17640000\) and
\(2\cdot4200\cdot37=310800\). These estimates preserve the already
calculated nonzero fourth-order kernel forms. They do not assert a new
lower bound on those forms. The uniform fourth-order identification and
a positive lower bound on the entire incidence kernel are proved below
in (16.21)--(16.24), using the retained coefficients and a smaller
explicit coupling interval.

There is also a uniform, quantitative local Rayleigh result. For any
single link \(e\) and \[
0<\xi\le\frac1{680000},                                      \tag{7.5}
\] (5.3), \(r_e\ge2\), and \(42500/680000=1/16\) give \[
C_{ee}\ge\frac{r_e}{32}\xi^2>0.
\] Equations (5.3) and (6.4) therefore imply \[
\left|\frac{\mathfrak q_A[v_e]}{\|v_e\|^2}-3\kappa\right|
 \le\frac{4944000}{r_e}\kappa\xi
 \le2472000\kappa\xi.                                       \tag{7.6}
\] The numerator in this calculation satisfies
\(|\mathcal N_{ee}-3\kappa C_{ee}|\le154500\kappa\xi^3\), which supplies
the constant in (7.6). This is the energy of the actual smooth nonzero
gauge-invariant local trial vector. By the spectral variational
principle its quotient is an upper bound on the physical finite-box
excitation gap, not a lower bound on that gap. The explicit range (7.5)
is equivalently \(g_{\rm YM}^4\ge170000\), retaining
\(\xi=1/(4g_{\rm YM}^4)\). Neither \(a\) nor \(L\) appears in this
range. Section 15 proves the matching volume-uniform lower threshold for
the physical operator and every nonzero signed electric state; together
they bound the attained minimum over the entire actual electric-state
family.

\section{8. Propagation to the original native magnetic
state}\label{propagation-to-the-original-native-magnetic-state}

For the unchanged \(n_2^2\) weights in (1.4), set \[
\mathcal A_L=\frac{L^2(2L+1)}{15}
 (24L^4+24L^3+8L^2-4L+3),
\] \[
S_{1,L}=\frac{2L^2(L+1)(2L+1)^2}{3},\qquad
S_{2,L}=\frac{2L^2(L+1)(2L+1)^2(3L^2+3L-1)}{15}.
\] The exact identities are
\(\|\mathsf Iw^{\rm nat}\|_2^2=4\mathcal A_L\),
\(\|w^{\rm nat}\|_1=S_{1,L}\), and \(\|w^{\rm nat}\|_2^2=S_{2,L}\). To
check them with all boundary terms, direction-1 links have \(2L\) first
coordinates and \(2L+1\) third coordinates, so the latter two sums are
\(2L(2L+1)\sum_{t=-L}^Lt^2\) and \(2L(2L+1)\sum_{t=-L}^Lt^4\). The exact
sums are \(L(L+1)(2L+1)/3\) and \(L(L+1)(2L+1)(3L^2+3L-1)/15\),
respectively; their formulas follow by taking the difference under
\(L\mapsto L+1\) and checking \(L=0\). For face sums the \(12\) profile
is \(t^2+(t+1)^2\), the \(13\) profile is \(2t^2\), and the \(23\)
profile is zero. Counting their allowed lower corners gives \[
\mathcal A_L=2L(2L+1)\sum_{t=-L}^{L-1}(t^2+t+1/2)^2
                  +4L^2\sum_{t=-L}^{L}t^4.
\] The first sum is \(L(4L^4+1)/10\), checked by the two added endpoint
terms and its value at \(L=1\); substitution gives the displayed
polynomial.

Equations (7.1)--(7.2) now yield the explicit new remainders \[
\left|\|v_{w^{\rm nat}}\|^2-\frac{\xi^2}{4}\mathcal A_L\right|
 \le42500\xi^3 S_{1,L}^2,                                    \tag{8.1}
\] \[
\left|\mathfrak q_A[v_{w^{\rm nat}}]
          -\frac{3\kappa\xi^2}{4}\mathcal A_L\right|
 \le351000\kappa\xi^3 S_{2,L}.                               \tag{8.2}
\] All native weights are nonnegative, so on each face
\((\sum w_e)^2\ge\sum w_e^2\); summing faces gives
\(4\mathcal A_L\ge\sum_e r_e(w_e^{\rm nat})^2\ge2S_{2,L}\). Consequently
the relative error in the native energy numerator is at most
\(936000\xi\), independently of \(L\). In particular on
\(0<\xi\le1/1872000\) it is at most one half of that numerator's
positive leading term. This is a proved fixed-coupling extensive energy
estimate, not a volume-shrinking perturbative statement.

The upper error-to-leading-term ratio furnished by (8.1) alone is
\(170000\xi S_{1,L}^2/\mathcal A_L\). Its \(L^3\) growth is visible in
the exact polynomials. This bound does not give a uniform relative error
for the native denominator. No step through Section 8 proves summability
of the separated electric correlations, and the fourth-order leading
kernel is not discarded. The proved local covariance and weighted
numerator estimates therefore do not yet establish the fixed-coupling
minimum over every extensive weight or its continuum limit. This
statement records the scope of the inequalities actually proved; it is
not an assertion that the remaining correlations have no relation to the
native state. Sections 11--12 below prove a uniform full spectral lower
bound and an extensive covariance upper bound on an explicit smaller
nonempty interval. Sections 13--14 then prove a uniform two-sided
remainder by controlling the separated correlations. In particular
(14.8) replaces the \(L^3\) relative loss by the explicitly vanishing
\(8\epsilon_\xi\), and (14.7) applies to the native electric-state
quotient. These later estimates do not change the original coefficients
in (8.1)--(8.2). The gauge-equivariant calculation in Section 15
additionally proves the uniform signed-state minimum (15.21) and
improves the native extensive upper bound to (15.17); the stronger lower
threshold is not obtained by discarding a small native or signed
covariance denominator.

The full physical coefficient conversions are \[
\xi^2=\frac1{16g_{\rm YM}^8},\quad
\xi^3=\frac1{64g_{\rm YM}^{12}},\quad
\kappa\xi^2=\frac1{8g_{\rm YM}^6a},\quad
\kappa\xi^3=\frac1{32g_{\rm YM}^{10}a}.
\] Thus, for example, the right side of (8.2) is exactly
\(351000S_{2,L}/(32g_{\rm YM}^{10}a)\). The actual vacuum energy
includes \(2bM=M/(g_{\rm YM}^2a)\) throughout. The proof does not
exchange any spatial, coupling, cusp, or continuum limits.

\section{9. Primary literature, exact dictionaries, and
provenance}\label{primary-literature-exact-dictionaries-and-provenance}

The local projection method is related to the Feshbach--Schur
reconstruction map, not claimed as a new general projection principle.
Dusson--Sigal--Stamm, \emph{The Feshbach--Schur map and perturbation
theory}, arXiv:2105.02058v1, Theorem 1.2, equation (1.11), and Lemma 2.1
give the complement inverse and eigenvector reconstruction mechanism
(\href{https://arxiv.org/abs/2105.02058}{primary source}). Their
finite-rank perturbation theorem is not used as a volume-uniform
hypothesis. Here \(P=P_S\) has infinite rank when the exterior is
nonempty. For an exact operator dictionary, set \[
A_S^Q=Q_S(H-\mathcal E)Q_S|_{Q_S\mathcal H_L}
 =B_S|_{Q_S\mathcal H_L}-bQ_SW_SQ_S.
\] On its domain, when \(2m_S\xi<3/4\), the bound
\(A_S^Q\ge\kappa(3/4-2m_S\xi)>0\) proves its bounded inverse. The
\(Q_S\) block of the full eigen-equation gives exactly \[
Q_S\psi=b(A_S^Q)^{-1}Q_SW_SP_S\psi.                            \tag{9.1}
\] The retained block is then \[
\left[K_S-b^2P_SW_SQ_S(A_S^Q)^{-1}Q_SW_SP_S\right]P_S\psi=0.   \tag{9.2}
\] Equations (9.1)--(9.2) follow directly by substitution, including the
minus sign and complete exterior operator. Their domains are the
exterior domain of \(K_S\) and its image under the bounded
reconstruction correction. The off-diagonal factors are bounded, so the
effective correction in (9.2) is bounded on the exterior Hilbert space.
Conversely a vector in that domain solving (9.2), inserted into (9.1),
lies in the full domain and solves the original zero-eigenvalue
equation. Applying \(P_S\) to this reconstruction recovers the original
exterior vector; its kernel is therefore zero. This proves the exact
bijection of the two kernels. The estimates in Sections 3--8 instead use
(2.8), whose inverse needs no smallness condition on \(b\); (9.1)
explains the precise relation between the two complement-resolvent
presentations.

The canonical corpus also routes Henheik--Teufel--Wessel, \emph{Local
stability of ground states in locally gapped and weakly interacting
quantum spin systems}, arXiv:2106.13780v3, Section 2 and Theorem 7
(\href{https://arxiv.org/abs/2106.13780}{primary source}). It permits
infinite-dimensional single-site spaces and unbounded on-site operators.
Its locality constants are existential in the statements read; none is
assigned a numerical value in this note. The exact model dictionary can
be given without changing physical coordinates. Associate the edge
\((n,i)\) to the indexing site \(x(e)=2n+\mathbf e_i\in\mathbb Z^3\),
retaining this map's inverse on its image: the unique odd coordinate
identifies \(i\), and then \(n=(x-\mathbf e_i)/2\). Attach
\(L^2(SU(2))\) and on-site \(h_{x(e)}=\kappa E_e\) to that site. The
on-site gap is \(3\kappa/4\). Assign a face \(p=(n;i,j)\), \(i<j\), to
its edge \((n,i)\). Its other three indexing sites have \(\ell^1\)
distance \(2\) from \(x(n,i)\). At most two faces are assigned to one
edge, so \(\Phi_{x(e)}=b\sum_{p\ {\rm assigned\ to}\ e}(2-W_p)\) has
norm at most \(8b\) and range \(2\). The tensor-permutation unitary
defined on elementary tensor functions by this bijection intertwines the
complete Hamiltonian with \(\sum_xh_x+\sum_x\Phi_x\), including all
scalar terms. The indexing map is not a physical rescaling: the original
midpoint is exactly \(a\,x(e)/2\). The source's small-interaction
theorem is not invoked to supply the explicit constants (3.1)--(8.2);
those were proved above for this precise model.

The local source identities and conventions were read in full in:

\begin{itemize}
\tightlist
\item
  \texttt{magnetic\_translation\_true\_vacuum.md}, Sections 1--9,
  especially (1.1), (1.4), (5.3), and (8.1)--(8.5);
\item
  \texttt{local\_energy\_current\_true\_vacuum.md}, Sections 1--6;
\item
  \texttt{wilson\_variation/ELECTRIC\_COVARIANCE\_UNSIGNED\_INCIDENCE\_20260908.md},
  Sections 1--8, especially (2.8)--(2.9), (5.7), (6.5), and (7.13);
\item
  \texttt{wilson\_variation/BULK\_PLAQUETTE\_EXTENSION.md}, Sections
  1--10, especially the explicit box-dependent disk and conditional-form
  maps.
\end{itemize}

Their common root is \texttt{../ym\_gap\_primary\_20260908}. The
complete finite incidence-kernel and second-order coefficient
calculations are preserved there. This note does not repeat those
calculations or promote their \(O_L\) constants to uniform ones. Its new
input to that chain is the independently proved quantitative
local-resolvent control and the resulting estimates above. The
machine-readable lane route records the canonical query results, primary
TeX locators, exact reading coverage, and remaining source reads.

\section{10. A uniform score bound and coercivity of the actual
conditional
operator}\label{a-uniform-score-bound-and-coercivity-of-the-actual-conditional-operator}

The full vacuum equation also controls the conditional density without a
small-coupling expansion. This gives a second volume-independent
estimate useful in the exact elimination construction. It does not
presuppose a spectral gap for the full interacting system.

Use the product Riemannian metric for which the three \(T_a\) fields at
each link are orthonormal; write \(\nabla_e,\Delta_e\) for its link
gradient and Laplacian, and \(\Delta=\sum_e\Delta_e=-\sum_eE_e\). This
is precisely the metric already used in (1.2), not a new physical
metric. In quaternion coordinates \(U=u_0I+i\sum_a u_a\sigma_a\), its
metric is four times the metric of the unit three-sphere: the tangent
vector \(T_a\) at the identity has quaternion Euclidean length \(1/2\),
while its specified Riemannian length is 1. Left multiplication is
orthogonal on quaternions, so this identity holds at every point.
Consequently a link factor is isometric to the round sphere of radius 2.
Its sectional curvature is \(1/4\), Ricci tensor is \(g/2\), and
diameter is \(2\pi\). The curvature coefficient also follows from the
second fundamental form of the radius-2 sphere: its normal derivative
has magnitude \(1/2\), the Gauss equation gives sectional curvature
\(1/4\), and summing the two transverse sectional curvatures gives
\(1/2\). These computations retain the factor four relative to \(c\).

Put \(u=\log\psi\), a globally smooth real function because \(\psi>0\).
Dividing the eigen-equation by the same positive \(\psi\), with its
original unit norm left unchanged, yields the exact nonlinear equation
\[
\Delta u+|\nabla u|^2=\frac{V-\mathcal E}{\kappa},
\qquad V=b\sum_{p\in\mathsf P_L}(2-W_p).                         \tag{10.1}
\] For each fixed link \(e\) define \(w_e=|\nabla_eu|^2\), and let
\(\nabla_f\nabla_eu\) denote the covariant derivative, in factor \(f\),
of the factor-\(e\) gradient. On the product manifold the partial
Bochner identity is \[
\Delta w_e
=2\sum_f|\nabla_f\nabla_eu|_{\rm HS}^2
 +2\langle\nabla_eu,\nabla_e\Delta u\rangle+w_e.                  \tag{10.2}
\] Here is its coordinate derivation. Take orthonormal geodesic frames
in every factor at the point under consideration. Applying each second
derivative in \(\Delta_f\) to the sum of squares of the three components
of \(\nabla_eu\) gives twice their squared first derivatives plus twice
their contraction with the second derivative of the vector. For
\(f\ne e\), the product connection has zero mixed curvature, so that
second derivative commutes with \(\nabla_e\). For \(f=e\), commuting the
covariant derivatives of the differential \(du\) gives \[
\Delta_e^{\rm rough}\nabla_eu
=\nabla_e\Delta_eu+\operatorname{Ric}_e(\nabla_eu).
\] This sign can equally be read from the positive sectional-curvature
Gauss formula above. The Ricci contraction in the differentiated norm is
\(2\operatorname{Ric}_e(\nabla_eu,\nabla_eu)=w_e\). Adding the factor
equations proves (10.2). The scalar smooth Bochner identity, with the
same sign of \(\Delta\), is given in C. Villani, \emph{Optimal
Transport: Old and New}, Chapter 14, equation (14.28), printed p.388
(PDF p.394 in the indexed local edition). The partial-factor and vacuum
substitutions used here are calculated explicitly, rather than inferred
from a dimension-independent constant not stated there.

Let \(\mathcal D_u=\Delta+2\langle\nabla u,\nabla\,\cdot\,\rangle\).
Differentiating (10.1) and using the symmetry of the full Hessian gives
\[
\mathcal D_u w_e
=2\sum_f|\nabla_f\nabla_eu|_{\rm HS}^2+w_e
 +\frac2\kappa\langle\nabla_eu,\nabla_eV\rangle.                 \tag{10.3}
\] In detail, the derivative of \(|\nabla u|^2\) in factor \(e\) is
\(2(\nabla^2u\,\nabla u)_e\). Its contraction with \(2\nabla_eu\) in
(10.2) is minus \(4\nabla^2u(\nabla u,\nabla_eu)\). The drift applied to
\(w_e\) is plus that same quantity. This proves the cancellation with
every mixed factor \(f\) included.

\textbf{Theorem 10.1.} For all \(L\ge2\), \(a,g_{\rm YM}>0\), and each
link, \[
\|\nabla_e\log\psi\|_\infty\le2r_e\xi,\qquad
\|\nabla_e\log\rho\|_\infty\le4r_e\xi,\qquad \rho=\psi^2.        \tag{10.4}
\] In particular \[
0\le\gamma_e\le4r_e^2\xi^2.                                   \tag{10.5}
\] This improves the constant in the earlier valid bound (3.1). All
estimates (4.1)--(8.2) remain valid with their displayed constants.

\textbf{Proof.} Equation (6.2) gives \(|\nabla_eW_p|\le1\) for every
face at \(e\), and all other face derivatives vanish. Thus
\(|\nabla_eV|\le br_e\), without any factor \(M\). The smooth
nonnegative function \(w_e\) attains its maximum on the compact product.
At that point its full gradient vanishes and \(\Delta w_e\le0\).
Equation (10.3) therefore gives \[
0\ge w_e-2r_e\xi\sqrt{w_e}.
\] If the maximum is zero the conclusion is immediate. Otherwise
dividing by its positive square root gives \(\sqrt{w_e}\le2r_e\xi\).
This proves the first bound everywhere. Since \(\log\rho=2u\), the
second follows with its exact factor two. Finally integration by parts
gives \(\gamma_e=\int|\nabla_e\psi|^2dU=\int\rho|\nabla_eu|^2dU\); its
probability integral and (10.4) prove (10.5). \(\square\)

To apply this to the actual conditional operator, fix a nonempty set
\(S\) of eliminated links and write \(r=U_{S^c}\), \(s=U_S\). Define \[
\rho_{S^c}(r)=\int\rho(s,r)\,ds,\qquad
p_S(s\mid r)=\frac{\rho(s,r)}{\rho_{S^c}(r)},\qquad
\omega_S=8\pi\xi\sum_{e\in S}r_e.                              \tag{10.6}
\] These are the marginal and conditional density of the same vacuum,
not a postulated Wilson Gibbs density. Holding \(r\) fixed and joining
two \(S\)-configurations one link at a time by minimizing geodesics,
(10.4) and the diameter \(2\pi\) give \[
\sup_s\log p_S(s\mid r)-\inf_s\log p_S(s\mid r)\le\omega_S.      \tag{10.7}
\] The denominator in (10.6) cancels from this oscillation. No
differentiability or uniform limit of the exterior marginal is assumed.

The product Haar Poincare inequality on \(S\), with its exact
single-link spectral coefficient, is \[
\inf_{c\in\mathbb C}\int|F-c|^2ds
\le\frac43\sum_{e\in S}\int|\nabla_eF|^2ds.                     \tag{10.8}
\] To prove it, expand \(F\) in the complete product matrix-coefficient
basis used in Section 2. Removing the constant leaves at least one
nonzero spin, so its total electric eigenvalue is at least \(3/4\).
Parseval gives (10.8) first for finite sums and then by form closure for
every \(H^1\) function. For each exterior \(r\), let
\(p_{\max},p_{\min}\) be the maximum and minimum of its smooth positive
conditional density on \(S\). Minimize over \(c\), first bound the
density above, apply (10.8), and then bound it below in the energy
integral. This gives the full comparison calculation \[
\begin{split}
\operatorname{Var}_{p_S(\cdot\mid r)}(F)
&\le p_{\max}\inf_c\int|F-c|^2ds\\
&\le\frac{4p_{\max}}3\sum_{e\in S}\int|\nabla_eF|^2ds\\
&\le\frac43 e^{\omega_S}
       \sum_{e\in S}\int|\nabla_eF|^2p_S(s\mid r)\,ds .
\end{split}                                                   \tag{10.9}
\] The minimum characterization of variance is valid for complex
functions as well, with the complex mean as minimizer. This proof
records the entire density ratio; it is the bounded-density perturbation
argument, not a replacement of either measure.

Let \(\mathscr H_\rho=L^2(\rho\,dU)\), and define the exact isometry \[
J_S:L^2(\rho_{S^c}\,dr)\longrightarrow\mathscr H_\rho,\qquad
(J_SF)(s,r)=F(r).
\] Its adjoint is \[
(J_S^*G)(r)=\int G(s,r)p_S(s\mid r)\,ds.
\] Fubini verifies the inner-product identity for the adjoint and
\(J_S^*J_S=I\). Thus \(P_S^\rho=J_SJ_S^*\) and \(Q_S^\rho=I-P_S^\rho\)
are orthogonal projections; \(\mathscr K_S^\rho=\ker J_S^*\) is the
conditional-mean-zero space. The multiplication map
\(U_\psi:G\mapsto\psi G\) is a unitary
\(\mathscr H_\rho\to\mathcal H_L\). The exact ground-state transform and
its closed form are \[
\mathscr L_\rho=U_\psi^{-1}(H-\mathcal E)U_\psi
 =-\kappa(\Delta+\langle\nabla\log\rho,\nabla\,\cdot\,\rangle),
\qquad
q_\rho[G]=\kappa\sum_e\int|\nabla_eG|^2\rho\,dU.                 \tag{10.10}
\] At each fixed box the density and its reciprocal are smooth and
bounded on a compact space. Differentiating the conditional integral
shows that \(P_S^\rho,Q_S^\rho\) preserve \(H^1\): a retained derivative
gives the conditional integral of the input derivative and a bounded
smooth derivative of the conditional kernel; an eliminated derivative of
the output is zero. The constants in this domain-preservation statement
need not be uniform, and are not used for the coercivity bound below.
Applying \(Q_S^\rho\) to smooth \(H^1\) approximants proves that
\(H^1\cap\mathscr K_S^\rho\) is the form closure of its smooth elements
and dense in \(\mathscr K_S^\rho\).

Define \(C_S^\rho\) as the self-adjoint operator of the restricted
closed form \(q_\rho|_{H^1\cap\mathscr K_S^\rho}\). Explicitly its
domain consists of those \(G\) in that form space for which there is
\(h\in\mathscr K_S^\rho\) satisfying
\(q_\rho(F,G)=\langle F,h\rangle_\rho\) for every \(F\) in the form
space, and \(C_S^\rho G=h\). This definition makes no unproved claim
about the domain of a formal operator product.

\textbf{Theorem 10.2.} For every positive coupling and every finite box,
\[
C_S^\rho\ge\frac{3\kappa}{4}e^{-\omega_S}I,\qquad
\|(C_S^\rho)^{-1}\|\le\frac4{3\kappa}e^{\omega_S}.              \tag{10.11}
\] In particular for a single eliminated link \(e\), \[
C_{\{e\}}^\rho\ge\frac{3\kappa}{4}e^{-8\pi r_e\xi}I
\ge\frac{3\kappa}{4}e^{-32\pi\xi}I.                            \tag{10.12}
\] These single-link bounds are independent of all exterior volume and
valid for the actual interacting conditional density.

\textbf{Proof.} A vector \(G\in\mathscr K_S^\rho\) has zero conditional
mean for almost every \(r\). Apply (10.9) to its conditional slices and
multiply by \(\rho_{S^c}(r)\); integrate in \(r\). Fubini gives \[
\|G\|_\rho^2\le\frac43 e^{\omega_S}
       \sum_{e\in S}\int|\nabla_eG|^2\rho\,dU
\le\frac4{3\kappa}e^{\omega_S}q_\rho[G].
\] Smooth approximation extends this inequality to the entire restricted
form domain. The form representation and spectral theorem yield (10.11);
\(r_e\le4\) gives (10.12). The unitary \(U_\psi\) carries this form onto
the full physical excitation form restricted to
\(U_\psi\mathscr K_S^\rho\), so the constant has not been transferred
between different operators without a domain map. Gauge transformations
preserve the density and intertwine conditional integration by endpoint
Haar invariance. Consequently these projections, forms, and inverses
preserve the physical gauge-invariant subspaces as well. \(\square\)

For comparison with the original parameters, the last bound in (10.12)
is \[
\frac{3g_{\rm YM}^2}{2a}
       \exp\!\left(-\frac{8\pi}{g_{\rm YM}^4}\right).
\] For larger eliminated sets the bound is also explicit, but its
dependence on \(\sum_{e\in S}r_e\) is retained. No volume-independent
extensive-block bound is inferred by deleting that sum. Nor does a
uniform conditional one-link gap alone tensorize into a uniform gap for
the correlated full density. The proof has established an inverse for
the actual eliminated operator, not a separated-correlation estimate or
the Millennium spectral conclusion. The full finite-box spectral bound
proved in Section 12 uses an independently convergent cluster
construction, not an assumed tensorization of these conditional
inequalities.

The indexed Villani PDF has stable unit ID
\nolinkurl{PUBUNIT-3B56736161C0B49850EA7433}, SHA-256
\nolinkurl{3AE2DB61CBFFADFEEF0DE738742F4F3CE81B21874DACB1AD2991A7000434675D}.
Its contents and the complete printed pp.387--389, including equation
(14.28), were read directly; the book is cited, not copied into this
repository. The page reference is to that exact local edition rather
than an assumed page number in every edition.

\section{11. An explicit convergent cluster coordinate construction for
the full
vacuum}\label{an-explicit-convergent-cluster-coordinate-construction-for-the-full-vacuum}

The source method used here is the commuting creation-coordinate
construction in D. A. Yarotsky, \emph{Quasi-particles in weak
perturbations of non-interacting quantum lattice systems},
arXiv:math-ph/0411042v1, Section 2, equations labelled mmo, c1, vsum,
thl and sumj in the retained source TeX
(\href{https://arxiv.org/abs/math-ph/0411042}{primary source}). That
section explicitly permits infinite-dimensional single-site spaces and
unbounded on-site operators. Its smallness constants are not numerical.
The following proof derives its own constants for the complete four-link
Wilson interaction. The method is not claimed new. No finite-spin
truncation is made, and no rescaling of the original Hamiltonian or
vacuum is made.

Write \[
H_0=\kappa\sum_eE_e,\qquad V_p=-bW_p,\qquad
H=H_0+2bM+\sum_pV_p,\qquad \gamma=\frac{3\kappa}{4}.
                                                               \tag{11.1}
\] Thus \(\|V_p\|\le J:=2b\), each face has \(s=4\) distinct links, and
each link belongs to at most \(d=4\) faces. The scalar \(2bM\) remains
in (11.1) and in the eigenvalue equation below. It commutes with every
coordinate transformation and has zero nonconstant component; these
facts, not deletion of its value, account for its absence from the
nonconstant fixed-point equation.

Let \(\Omega=1\) in the original product Haar probability space. For
\(I\subset\mathsf E_L\), put \[
\mathcal H'_I=\bigotimes_{e\in I}
       \bigl(L^2(SU(2))\ominus\mathbb C1\bigr),\qquad
\mathcal H'_\varnothing=\mathbb C.
\] The canonical orthogonal tensor decomposition gives a unitary map \[
\mathcal T:\bigoplus_{I\subset\mathsf E_L}\mathcal H'_I
 \longrightarrow\mathcal H_L,\qquad
\mathcal T(z)=\sum_I z_I\otimes1_{I^c}.                         \tag{11.2}
\] Its inverse component \(\Pi_I\) is obtained by Haar integration over
\(I^c\) and application of the nonconstant projection in every link of
\(I\). Product projections prove both inverse identities in (11.2). For
nonempty \(I\), set \(H_I=\kappa\sum_{e\in I}E_e\) on \(\mathcal H'_I\).
The separate nonzero spins prove \[
H_I\ge\gamma |I|,\qquad
\|H_I^{-1}\|\le(\gamma |I|)^{-1}.                              \tag{11.3}
\]

For \(z_I\in\mathcal H'_I\), define the bounded operator, on the \(I\)
factor and tensored with the exterior identity, by \[
\widehat z_I f=z_I\langle1_I,f\rangle_{L^2(SU(2)^I)}.
                                                               \tag{11.4}
\] The inner product is the one fixed in Section 1, so (11.4) is linear
in \(f\). Its norm is \(\|z_I\|\). Integration of any one common
nonconstant link proves \[
\widehat z_I\widehat y_K=
\begin{cases}
0,&I\cap K\ne\varnothing,\\
\widehat{z_I\otimes y_K},&I\cap K=\varnothing.
\end{cases}                                                   \tag{11.5}
\] For entangled vectors the same calculation follows by approximation
by finite elementary tensors in \(\mathcal H'_I,\mathcal H'_K\);
operator norms converge by (11.4). In particular all these operators
commute. A product of more than \(N=|\mathsf E_L|\) nonempty creation
operators is zero.

Use the graph whose vertices are physical links and whose distinct
vertices are adjacent precisely when they share a face. For nonempty
\(I\), let \(\ell(I)\) be the smallest number of edges of a connected
subgraph containing \(I\); extra vertices are allowed. The finite
full-box graph is connected. Define \[
\omega(I)=2^{\ell(I)+1}.
\] Monotonicity under inclusion holds. For a face support
\(p=\partial p\) the induced graph is a four-vertex clique, so
\(\ell(p)=3\) and \(\omega(p)=16\). If each \(I_j\) meets \(p\),
connected minimizing subgraphs and a tree on \(p\) have connected union.
Therefore \[
\omega\left(p\cup\bigcup_{j=1}^n I_j\right)
 \le16\prod_{j=1}^n\omega(I_j).                               \tag{11.6}
\] No replacement of physical coordinates is involved in this graph.

Take the Banach space of collections \(z=(z_I)_{I\ne\varnothing}\),
\(z_I\in\operatorname{Dom}(H_I)\), with norm \[
\|z\|_*=\max_x\sum_{I\ni x}
     \omega(I)\frac{\|H_Iz_I\|}{\gamma}.
                                                               \tag{11.7}
\] There are finitely many \(I\), and each energy norm is complete by
the bounded inverse in (11.3), so (11.7) is complete. In particular,
writing \(a_I=\omega(I)\|z_I\|\), one has \[
\sup_x\sum_{I\ni x}|I|a_I\le\|z\|_*,
\qquad
\sum_{I:I\cap p\ne\varnothing}a_I\le s\|z\|_* .                \tag{11.8}
\] Let \(C(z)=\sum_{I\ne\varnothing}\widehat z_I\). It is bounded in
each finite box. Its exponential and inverse are the finite polynomials
\(e^{C(z)}\) and \(e^{-C(z)}\) by (11.5). On the full domain of \(H_0\),
\[
[H_0,\widehat z_I]=\widehat{H_Iz_I},\qquad
e^{-C(z)}H_0e^{C(z)}
 =H_0+\sum_{I\ne\varnothing}\widehat{H_Iz_I}.                  \tag{11.9}
\] To check domains, the exterior part of \(H_0\) commutes with (11.4);
the local part annihilates the Haar bra, and its action on the ket is
\(H_Iz_I\). This first proves the commutator on finite spectral tensors.
Its bounded right-hand side extends it to the full graph domain by
approximation. Thus \(C(z)\), its powers, and both exponentials preserve
that domain. The second commutator in (11.9) vanishes by (11.5), proving
the exact conjugation formula.

The fixed-point map is \[
\mathcal F(z)_I=-H_I^{-1}\Pi_I
  e^{-C(z)}\left(\sum_pV_p\right)e^{C(z)}\Omega,
       \qquad I\ne\varnothing.                              \tag{11.10}
\] The projection output is an arbitrary vector of \(\mathcal H'_I\);
(11.3) places \(\mathcal F(z)_I\) in the correct operator domain. We now
bound (11.10), including all clusters, rather than assuming the
convergence estimate quoted in the source.

For a fixed face \(p\), expand \[
e^{-C(z)}V_pe^{C(z)}
 =\sum_{n\ge0}\frac1{n!}
    [\cdots[[V_p,C(z)],C(z)],\ldots,C(z)].
                                                               \tag{11.11}
\] Only indices \(I_j\) meeting \(p\) contribute: a creation operator
disjoint from \(p\) commutes with \(V_p\) and with all other creation
operators. Every \(n\)-fold nested commutator expands into \(2^n\)
ordered words. A nonzero word has disjoint creation supports on each
side of \(V_p\). Since they all meet the four-link support of \(V_p\),
there are at most four on either side. Thus (11.11) actually ends at
\(n=8\). Infinite exponential sums used below are upper bounds for this
finite exact polynomial.

There is a dimension-independent bound on the tensor-sector sum of each
word acting on \(\Omega\). Outside \(p\), \(V_p\) acts as identity. A
common outside link of two creation supports annihilates the word, by
integrating a nonconstant factor. Otherwise the outside excitation set
is exactly \((\bigcup_jI_j)\setminus p\). Only the choice of excited
links inside \(p\) varies in the output. There are at most \(2^s=16\)
such orthogonal sectors. Cauchy--Schwarz and (11.4) give \[
\sum_K\|\Pi_K(\hbox{word})\Omega\|
 \le2^{s/2}J\prod_{j=1}^n\|z_{I_j}\|
 =4J\prod_{j=1}^n\|z_{I_j}\|.                                \tag{11.12}
\] The argument is unchanged for entangled creation vectors by the
elementary-tensor approximation already used in (11.5).

Each nonempty output sector \(K\) is contained in \(p\cup\bigcup_jI_j\).
To bound the anchored sum for \(x\in K\), cover its possibilities by
\(x\in p\) and \(x\in I_j\). For the first possibility there are at most
\(d\) choices of \(p\); (11.8) bounds each remaining index sum by
\(sR\), where \(R=\|z\|_*\). For the second, choose \(j\), then
\(I_j\ni x\). At most \(d|I_j|\) faces meet that set; the first
inequality in (11.8) absorbs this full factor \(|I_j|\). Thus the
respective anchored bounds for the product sums are \[
d(sR)^n,\qquad ndR(sR)^{n-1}.                                \tag{11.13}
\] The second quantity is zero for \(n=0\). Combining (11.6),
(11.10)--(11.13), and the exact cancellation of \(H_I\) against its
inverse yields \[
\|\mathcal F(z)\|_*
\le\frac{4Jd\,16}{\gamma}
 \sum_{n\ge0}\frac{2^n}{n!}
       \bigl[(sR)^n+nR(sR)^{n-1}\bigr]
=\frac{2048}{3}\xi(1+2R)e^{8R}.                              \tag{11.14}
\] This anchor counting, rather than the norm of the full perturbation,
removes the factor \(M\).

For two collections in a ball of radius \(R\), telescope each
multilinear word in (11.11), inserting their difference in one position
at a time. The counting (11.13) holds with the norm of that difference
in its position, whether the anchor is there or elsewhere. For the
homogeneous degree-\(n\) bound this multiplies its coefficient by
\(nR^{n-1}\) instead of \(R^n\). Differentiating the positive majorant
series in (11.14) proves \[
\|\mathcal F(z)-\mathcal F(y)\|_*
 \le\frac{2048}{3}\xi(10+16R)e^{8R}\|z-y\|_* .                \tag{11.15}
\] This is a bound on the exact finite commutator map, not on a
truncated approximation to the vacuum.

Set \(R_*=1/16\). Since
\(\sum_{n\ge0}(1/2)^n/n!<\sum_{n\ge0}(1/2)^n=2\), (11.14)--(11.15) give,
on the explicit nonempty interval \[
0<\xi\le\frac1{49152},                                      \tag{11.16}
\] \[
\|\mathcal F(z)\|_*\le1536\xi\le\frac1{32}<R_*,
\qquad
\operatorname{Lip}(\mathcal F)\le\frac{45056}{3}\xi
 \le\frac{11}{36}<1.                                       \tag{11.17}
\] Iteration from zero therefore converges in the complete ball to a
unique fixed point \(z\). More explicitly, successive increments are
bounded by a geometric series of ratio \(11/36\); the limit solves
(11.10) by continuity, and applying (11.15) to two fixed points proves
uniqueness. Complex conjugation preserves the real Wilson potential, the
Haar projections and the inverses \(H_I^{-1}\). It commutes with the
iteration, so \(z\), \(e^{C(z)}\Omega\), and the eigenvalue below are
real.

Equations (11.9)--(11.10) show that \[
\phi=e^{C(z)}\Omega,\qquad
H\phi=\varepsilon\phi,\qquad
\varepsilon=2bM+
 \Pi_\varnothing e^{-C(z)}\left(\sum_pV_p\right)e^{C(z)}\Omega.
                                                               \tag{11.18}
\] Every nonempty creation product has zero Haar integral, hence
\(\langle\Omega,\phi\rangle=1\), so \(\phi\ne0\). The next section
proves that \(\varepsilon=\mathcal E\), and identifies the exact scalar
relating \(\phi\) to the unchanged physical unit vacuum. No equality of
these differently specified vectors is presumed.

\section{12. Uniform full spectral estimate and the extensive electric
covariance}\label{uniform-full-spectral-estimate-and-the-extensive-electric-covariance}

In this section (11.16) holds. Fix its constructed \(z\) and abbreviate
\(C=C(z)\). The bounded invertible similarity \[
\widetilde H=e^{-C}(H-\varepsilon)e^C=H_0+\mathcal K,\qquad
\widetilde H\Omega=0                                       \tag{12.1}
\] has the full domain of \(H_0\), by (11.9). It is not asserted
unitary. Both similarities and the exact physical Hamiltonian, including
its ground energy, remain explicit.

Give the finite orthogonal sector decomposition the additional norm \[
\|f\|_{\oplus,1}=\sum_I\|\Pi_If\|,\qquad
\|f\|\le\|f\|_{\oplus,1}\le2^{N/2}\|f\|.                     \tag{12.2}
\] The second inequality is Cauchy--Schwarz over all \(2^N\) sectors.
Thus the norm has exactly the same vectors and topology in each fixed
box. The equivalence constant depends on volume; no uniform claim about
that equivalence is used.

For \(u_I\in\operatorname{Dom}(H_I)\), (11.5), (11.9), (12.1) and
\(\widetilde H\Omega=0\) give exactly \[
\mathcal K\,\widehat u_I\Omega
=e^{-C}\left[\sum_pV_p,\widehat u_I\right]e^C\Omega .
                                                               \tag{12.3}
\] In detail, commute \(\widetilde H\) with \(\widehat u_I\); the
\(H_0\) commutator supplies \(H_Iu_I\), while the creation sum in (11.9)
commutes with \(\widehat u_I\). For the remaining term, \(e^C\) commutes
with \(\widehat u_I\), proving (12.3). Faces disjoint from \(I\) have
zero commutator, so at most \(d|I|\) faces remain.

Expand each remaining conjugation as in (11.11). All additional creation
indices still have to meet \(p\): the derivation identity and
commutation with \(\widehat u_I\) allow the nested commutators to be
written as the commutator with \(\widehat u_I\) of the corresponding
conjugated \(V_p\). The output outside \(p\) now has the fixed excited
support \((I\cup\bigcup_jI_j)\setminus p\), or the word vanishes. Thus
(11.12) applies with an extra factor \(\|u_I\|\). The commutator with
\(\widehat u_I\) contributes two words. Using
\(\sum_{K:K\cap p\ne\varnothing}\|z_K\|\le sR_*\) gives \[
\|\mathcal K\widehat u_I\Omega\|_{\oplus,1}
\le 8Jd\,|I|\,e^{8R_*}\|u_I\|.
                                                               \tag{12.4}
\] The bound also applies to the empty component with value zero.
Summing (12.4) and using (11.3) proves the relative operator bound \[
\|\mathcal Kf\|_{\oplus,1}
\le c_\xi\|H_0f\|_{\oplus,1},\qquad
c_\xi=\frac{512}{3}\xi\le\frac1{288}<1,
       \quad f\in\operatorname{Dom}(H_0).                    \tag{12.5}
\] Indeed \(8Jd/\gamma=256\xi/3\), and \(e^{8R_*}=e^{1/2}<2\). The bound
on finite spectral tensors extends in the graph norm; all factors in
(12.3) preserve the required domain as proved above.

Here is the resolvent argument with no appeal to a missing uniform
perturbation theorem. On every nonempty sector \(I\),
\(\operatorname{Spec}(H_I)\subset[\gamma|I|,\infty)\). For real \(t<0\),
functional calculus gives \[
\|\mathcal K(H_0-t)^{-1}\|_{\oplus,1}
 \le c_\xi\sup_{\lambda\ge\gamma}
          \frac{\lambda}{\lambda-t}\le c_\xi<1.
\] For \(0<t<(1-c_\xi)\gamma\), it gives instead \[
\|\mathcal K(H_0-t)^{-1}\|_{\oplus,1}
 \le\frac{c_\xi\gamma}{\gamma-t}<1.                           \tag{12.6}
\] The empty sector causes no term because \(\mathcal K\Omega=0\). The
geometric Neumann series for \((I+\mathcal K(H_0-t)^{-1})^{-1}\),
followed by \((H_0-t)^{-1}\), is a bounded inverse of
\(\widetilde H-t\). The domain is verified by the resolvent's
graph-domain image. Norm equivalence (12.2) makes it also a bounded
inverse in the original Hilbert space, without changing the numerical
spectral exclusion (12.6). Similarity (12.1) transfers this inverse to
\(H-\varepsilon-t\).

The original \(H\) is self-adjoint, \(\varepsilon\) is its real
eigenvalue, and no spectrum lies below it by the \(t<0\) argument. Thus
\(\varepsilon=\mathcal E\). Its simple positive ground vector was proved
in Section 1. The exact identification is therefore \[
\psi=\alpha e^C\Omega,\qquad
\alpha=\langle\Omega,\psi\rangle>0,\qquad
\alpha^2\|e^C\Omega\|^2=1.                                  \tag{12.7}
\] This records the physical scalar and preserves the original unit
vacuum rather than replacing it with a coefficient-normalized vector.
The spectrum above it obeys the explicit volume-independent estimate \[
(H-\mathcal E)\big|_{\psi^\perp}
 \ge\gamma(1-c_\xi)I
 =\frac{3\kappa}{4}\left(1-\frac{512}{3}\xi\right)I
 \ge\frac{287\kappa}{384}I .                                \tag{12.8}
\] This follows from (12.6) and the spectral theorem, since the ground
eigenvalue is simple. Restricting the same inequality to the
gauge-invariant subspace retains it, because that closed subspace
reduces \(H\) and contains \(\psi\). Section 15 proves that the cluster
similarity itself intertwines this constraint. Its restricted free
threshold is exactly \(3\kappa\), so (15.13) improves the physical
restriction of (12.8), not the full-space inequality or the numerical
fixed-point constants. In original parameters its range and last lower
bound are \[
g_{\rm YM}^4\ge12288,\qquad
\frac{287g_{\rm YM}^2}{192a}.
                                                               \tag{12.9}
\] This is a strong-coupling, fixed-spatial-regulator theorem for the
full interacting finite boxes. No spatial continuum limit or Millennium
conclusion follows by deleting its coupling range.

The spectral estimate has an immediate quantitative consequence for the
previously unresolved extensive covariance upper bound. Every \(v_w\) in
Section 7 is perpendicular to \(\psi\), so the spectral theorem and
(12.8) imply \[
\|v_w\|^2\le
\frac{4}{3\kappa(1-c_\xi)}\,\mathfrak q_{H-\mathcal E}[v_w].
\] Combining this with the already proved all-weight estimate (7.2)
gives, for every real weight vector and every box, \[
0\le w^TCw\le
\frac{\xi^2\|\mathsf Iw\|_2^2/4
             +468000\xi^3\|w\|_2^2}{1-512\xi/3}.
                                                               \tag{12.10}
\] The norm of \(\mathsf I\) is at most 4: on each face,
\((\sum_{e\in\partial p}w_e)^2\le4\sum_{e\in\partial p}w_e^2\); summing
faces and retaining \(r_e\le4\) proves
\(\|\mathsf Iw\|_2^2\le16\|w\|_2^2\). Consequently the fully extensive
matrix upper bound is \[
0\le C\le
\frac{4\xi^2+468000\xi^3}{1-512\xi/3}\,I_{\mathbb R^{\mathsf E_L}}.
                                                               \tag{12.11}
\] This removes the \(\ell^1\)-squared volume loss from an upper bound
on the exact covariance. This argument alone does not assert the
entrywise absolute row summability of \(C\), or an \(\ell^2\) remainder
of order \(\xi^3\) around \(\xi^2\mathsf I^T\mathsf I/16\).
Positive-semidefinite order and absolute-value entry sums are not
interchanged. Sections 13--14 prove the separate absolute-row and
two-sided claims by an explicit locality and spectral-filter
calculation.

For the original native weights the exact substitution is \[
\|v_{w^{\rm nat}}\|^2\le
\frac{\xi^2\mathcal A_L+468000\xi^3S_{2,L}}{1-512\xi/3}.
                                                               \tag{12.12}
\] Using \(S_{2,L}\le2\mathcal A_L\), proved in Section 8, yields \[
\|v_{w^{\rm nat}}\|^2\le
\xi^2\mathcal A_L
\frac{1+936000\xi}{1-512\xi/3}.
                                                               \tag{12.13}
\] Unlike (8.1), this upper control has no \(L^3\) relative loss.
Equation (8.1) remains a valid two-sided remainder with its original
constants; (12.13) supplies a different, stronger extensive upper
control on the nonempty range (11.16). On the entire leading kernel
\(\mathsf Iw=0\), (12.10) also yields \[
\|v_w\|^2\le\frac{468000\xi^3}{1-512\xi/3}\|w\|_2^2.
                                                               \tag{12.14}
\] This complements, and does not discard, the fourth-order
\(\ell^1\)-weighted bound (7.3) and the parent's exact fourth-order
calculation. The upper bounds (12.10)--(12.14) alone do not give a
two-sided estimate for the minimum over all signed actual states. The
uniform minimum, as opposed to every individual signed quotient, is
obtained in (15.21) by the gauge-restricted spectral argument. Equations
(15.15)--(15.17) also strengthen these PSD upper bounds by a factor of
four. The original fourth-order kernel coefficients are retained;
(16.21) supplies their uniform covariance and energy remainders.

Finally the same result improves exact extensive elimination, not only
the electric variational family. For every nonempty eliminated set
\(S\), a vector \(G\in\ker J_S^*\) satisfies \[
\int G\rho\,dU=\int (J_S^*G)(r)\rho_{S^c}(r)\,dr=0.
\] Thus \(U_\psi G=\psi G\) is perpendicular to the actual ground
vector. Apply (12.8) to its form and use (10.10), first on smooth
vectors and then by the proved form closure. This yields \[
C_S^\rho\ge\frac{3\kappa}{4}(1-512\xi/3)I
 \ge\frac{287\kappa}{384}I
 \quad\left(0<\xi\le\frac1{49152}\right),                    \tag{12.15}
\] \[
\|(C_S^\rho)^{-1}\|
 \le\frac{4}{3\kappa(1-512\xi/3)}
 \le\frac{384}{287\kappa}.                                  \tag{12.16}
\] Every constant in (12.15)--(12.16) is independent of the size, shape
and location of \(S\), as well as the exterior volume. The operator is
precisely the restricted full form \(C_S^\rho\) defined in Section 10,
containing derivatives in all link directions. It is not replaced by a
conditional-fibre differential operator containing only eliminated
derivatives. The all-positive-coupling bound (10.11) remains available
outside (11.16). On their common range both lower bounds hold, so the
larger of the two constants may be used. The unitary and gauge maps
remain exactly those proved in Section 10.

For the reducing gauge-invariant conditional-complement space,
(15.22)--(15.23) prove the sharper inverse constant \(96/(287\kappa)\).
The projection and unitary supplying that restriction are explicit.

\section{13. Absolute decay of the unbounded electric
covariance}\label{absolute-decay-of-the-unbounded-electric-covariance}

Throughout Sections 13--14, \(0<\xi\le1/49152\). Retain \(H,\psi\) and
\(E_e\) exactly as in Section 1, and set \[
\Delta_*=\frac{287\kappa}{384},\qquad
Q_\psi=I-|\psi\rangle\langle\psi|,\qquad
\alpha_t(O)=e^{itH}Oe^{-itH}.
\] Equation (12.8), not an additional assumption, gives
\(A|_{\psi^\perp}\ge\Delta_*\), where \(A=H-\mathcal E\). Let \(d(e,f)\)
be distance in the shared-face graph of Section 11. We will prove \[
|C_{ef}|\le1200\xi e^{-d(e,f)/4}\qquad(e\ne f).               \tag{13.1}
\] This is a bound on the original electric operators, which are
unbounded. No finite-spin Hamiltonian or approximate ground state is
introduced.

The two primary methods are B. Nachtergaele and R. Sims, \emph{On the
dynamics of lattice systems with unbounded on-site terms in the
Hamiltonian}, arXiv:1410.8174v1, Sections 2--3, especially Proposition
2.1, Lemma 2.2 and the proof of Theorem 3.1
(\href{https://arxiv.org/abs/1410.8174v1}{source}), and the same
authors, \emph{Lieb-Robinson Bounds and the Exponential Clustering
Theorem}, Commun. Math. Phys. 265 (2006), 119--130,
arXiv:math-ph/0506030v3, Theorem 2, Section 3.2 and Lemma 1
(\href{https://arxiv.org/abs/math-ph/0506030v3}{source}). The former
explicitly repairs domain and strong-continuity issues in earlier
unbounded-on-site proofs. Both source TeX files are retained. We prove
the numerical bounds and the electric-operator passage below; neither is
obtained by assigning constants to an existential statement.

\subsection{13.1. A second electric moment and an exact cutoff
comparison}\label{a-second-electric-moment-and-an-exact-cutoff-comparison}

For \(S=\{e\}\), joint spectral calculus in (2.7) gives, for
\(f\in\operatorname{Dom}(E_e)\), \[
\|E_e^2R_{\{e\}}f\|\le\kappa^{-1}\|E_ef\|.
\] Indeed, on a nonconstant link sector the multiplier satisfies
\(\lambda_e^2/(\kappa\lambda_e+k)\le\lambda_e/\kappa\); it vanishes on
the constant link sector. Applying this to the smooth vector
\(W_e^\star\psi\) in (2.8), and noting \(E_e^2P_{\{e\}}\psi=0\), proves
\[
\|E_e^2\psi\|\le\xi\|E_e(W_e^\star\psi)\|
\le\frac{3r_e}{2}\xi+8r_e^2\xi^2
\le6\xi+128\xi^2\le7\xi.                                   \tag{13.2}
\] Here the full product rule is \[
E_e(W_e^\star\psi)=\tfrac34W_e^\star\psi+
 W_e^\star E_e\psi-2D_e\psi.
\] The three terms have bounds \(3r_e/2\), \(4r_e^2\xi\), and
\(4r_e^2\xi\), respectively: the last follows from
\(\|D_e\psi\|\le r_e\sqrt{\gamma_e}\le2r_e^2\xi\), using (10.5). Thus
the mixed derivative has been estimated, not omitted.

For any \(\Lambda>0\), define the bounded one-link observable \[
B_e^\Lambda=E_e\mathbf1_{[0,\Lambda]}(E_e),\qquad
\|B_e^\Lambda\|\le\Lambda.
\] It is a functional calculus of the original \(E_e\); it changes
neither \(H\) nor \(\psi\). For its spectral tail,
\(\lambda\mathbf1_{\lambda>\Lambda}\le\lambda^2/\Lambda\), hence \[
\|(E_e-B_e^\Lambda)\psi\|\le7\xi/\Lambda,\qquad
\|Q_\psi B_e^\Lambda\psi\|\le\|B_e^\Lambda\psi\|
 \le\|E_e\psi\|\le8\xi.                                     \tag{13.3}
\] Define \(\operatorname{Cov}_\psi(O,P)=
\langle Q_\psi O\psi,Q_\psi P\psi\rangle\) for self-adjoint observables
with \(\psi\) in their domains. Subtracting the two centered inner
products, retaining both difference terms, gives \[
|C_{ef}-\operatorname{Cov}_\psi(B_e^\Lambda,B_f^\Lambda)|
\le112\xi^2/\Lambda.                                        \tag{13.4}
\] Both original link Casimirs and their spectral projections commute
with endpoint gauge transformations; every cutoff observable is
physical.

\subsection{13.2. Locality for the complete
dynamics}\label{locality-for-the-complete-dynamics}

For bounded observables \(O_e,O_f\) on distinct links, we claim \[
\|[\alpha_t(O_e),O_f]\|
\le2\|O_e\|\|O_f\|e^{-d(e,f)}
       (e^{v|t|}-1),\qquad v=192b.                           \tag{13.5}
\] This intermediate statement holds at every positive coupling. Here
are the domain construction and the full counting proof.

For a finite link set \(X\), introduce the auxiliary self-adjoint
operator \[
H_X^0=H_0+2bM+\sum_{\partial p\subset X}V_p,\qquad V_p=-bW_p.
\] It has the same \(H^2\) domain as \(H\), and its dynamics preserves
the algebra of bounded operators supported on \(X\). Its difference from
\(H\) is bounded. The propagator \[
W_X(t,s)=e^{itH_X^0}e^{-i(t-s)H}e^{-isH_X^0}
\] has inverse \(W_X(s,t)\), and on the common domain differentiating
gives \[
\partial_tW_X(t,s)=-iK_X(t)W_X(t,s),\qquad
K_X(t)=e^{itH_X^0}(H-H_X^0)e^{-itH_X^0}.
\] This bounded self-adjoint generator is strongly continuous. Its
vectorwise Dyson integrals exist, and on each compact time interval
their norms are bounded by the exponential series with parameter
\(\sup_t\|K_X(t)\|\). Iteration of the integral equation gives the
series; iteration for the difference of two solutions gives zero, since
its \(n\)-th bound has a factorial denominator. Consequently the domain
derivative agrees with that strong solution on the whole Hilbert space.
This justifies the differentiation below for bounded observables without
assuming operator-norm continuity of the unbounded on-site evolution.

Let \(\beta_{X,t}(O)=W_X(0,t)OW_X(t,0)\), and let
\(\partial_{\mathsf P}X\) be the faces meeting both \(X\) and \(X^c\).
For \(O\) supported on \(X\), the strong commutator derivative and
Jacobi identity give \[
\begin{split}
F(t)&=[\beta_{X,t}(O),O_f],\\
F'(t)&=i[L_X(t),F(t)]
   +i[\beta_{X,t}(O),[O_f,L_X(t)]],\\
L_X(t)&=\sum_{p\in\partial_{\mathsf P}X}\alpha_t(V_p).
\end{split}
\] Terms outside \(X\) commute with \(O\), and terms inside \(X\) are
already in \(H_X^0\); the crossing faces are therefore exactly the
stated ones. Conjugate this equation by the unitary strong propagator
generated by \(-L_X(t)\). Integrating its remaining inhomogeneous term
and taking norms proves \[
\|F(t)\|\le\|[O,O_f]\|+
 2\|O\|\sum_{p\in\partial_{\mathsf P}X}
 \int_0^{|t|}\|[\alpha_{\operatorname{sgn}(t)s}(V_p),O_f]\|\,ds.
\] At any fixed final \(t\), insert \(O=e^{itH_X^0}O_Xe^{-itH_X^0}\),
which still has support \(X\) and norm \(\|O_X\|\). This is a pointwise
use of the already proved inequality, not differentiation of a
time-dependent initial observable. It yields the recurrence \[
\begin{split}
\|[\alpha_t(O_X),O_f]\|
&\le2\|O_X\|\|O_f\|\mathbf1_{f\in X}\\
&\quad+2\|O_X\|\sum_{p\in\partial_{\mathsf P}X}
 \int_0^{|t|}\|[\alpha_{\operatorname{sgn}(t)s}(V_p),O_f]\|\,ds .
\end{split}                                                   \tag{13.6}
\] The Hilbert spaces here are separable. Every strongly continuous
commutator has a measurable norm, as the supremum over a countable dense
set of unit vectors; these scalar integrals are well-defined.

Iterate (13.6) from \(X=\{e\}\). A term reaching \(f\) in \(n\)
interactions is a chain of \(n\) faces, the first containing \(e\),
consecutive faces intersecting, and the last containing \(f\). Such a
chain gives a shared-face link path of at most \(n\) steps; therefore
\(n\ge d(e,f)\). There are at most \(4\) choices of the first face and
\(16\) choices of every subsequent face, since a face contains four
links and each link meets at most four faces. Each interaction has norm
at most \(2b\). The remainder after \(n\) iterations is bounded by a
constant times \((64b|t|)^n/n!\) and tends to zero. Dropping only
counted positive terms gives \[
\|[\alpha_t(O_e),O_f]\|
\le2\|O_e\|\|O_f\|\sum_{n\ge d(e,f)}
 \frac{(64b|t|)^n}{n!}
\le2\|O_e\|\|O_f\|e^{-d(e,f)}
       (e^{64\exp(1)b|t|}-1).
\] The last inequality uses \(\exp(n-d)\ge1\) for \(n\ge d\). The series
for \(\exp(1)\), with \(n!\ge2^{n-1}\) for \(n\ge2\) and a strict
inequality for \(n\ge3\), gives \(\exp(1)<3\). This proves (13.5),
including its explicit \(192b\). The constant \(2bM\) has commuted
through every displayed conjugation; the original Hamiltonian and its
energy origin remain unchanged.

\subsection{13.3. Spectral filtering with the actual vacuum
moments}\label{spectral-filtering-with-the-actual-vacuum-moments}

Take bounded self-adjoint \(O_e,O_f\), set
\(q_e=Q_\psi O_e\psi,\ q_f=Q_\psi O_f\psi\), and put
\(s_0=\|q_e\|\|q_f\|\). For \[
h(t)=\langle\psi,[O_e,\alpha_t(O_f)]\psi\rangle
=\langle q_e,e^{itA}q_f\rangle-\langle q_f,e^{-itA}q_e\rangle
\] one has \(|h(t)|\le2s_0\). For \(\alpha>0\), introduce the scalar
integral \[
I_\alpha=\lim_{\epsilon\downarrow0}\frac1{2\pi}
 \int_{\mathbb R}\frac{e^{-\alpha t^2}h(t)}{\epsilon+it}\,dt.
\] The limit is absolutely dominated: (13.5) makes \(h(t)=O(|t|)\) at
zero, while the Gaussian controls infinity. For \(\epsilon>0\), use
\((\epsilon+it)^{-1}=\int_0^\infty e^{-\epsilon u-itu}\,du\) and
integrate the Gaussian in \(t\). Fubini is justified by its integrable
absolute majorant. Letting \(\epsilon\downarrow0\) gives the bounded
spectral multiplier \[
F_\alpha(\lambda)=\frac1{\sqrt{4\pi\alpha}}
 \int_0^\infty e^{-(u-\lambda)^2/(4\alpha)}\,du
\] and the exact identity \[
I_\alpha=\langle q_e,F_\alpha(A)q_f\rangle
          -\langle q_f,F_\alpha(-A)q_e\rangle.
\] The Gaussian transform itself follows by differentiating its integral
with respect to its Fourier variable, integrating by parts, and fixing
the constant at zero by the squared Gaussian integral. For
\(\lambda\ge\Delta_*\), substitution followed by \((x+y)^2\ge x^2+y^2\)
gives \[
0\le1-F_\alpha(\lambda)=F_\alpha(-\lambda)
\le\tfrac12e^{-\Delta_*^2/(4\alpha)}.
\] The total variation of a cross spectral measure is at most
\(\|q_e\|\|q_f\|\): on a disjoint measurable partition apply
Cauchy--Schwarz to the two sequences of spectral-projection norms. Thus
every spectral interchange above is controlled, and \[
|\operatorname{Cov}_\psi(O_e,O_f)-I_\alpha|
\le s_0e^{-\Delta_*^2/(4\alpha)}.
\] Splitting its time integral at \(T>0\), (13.5) controls the part
\(|t|\le T\), while \(2s_0\) controls the rest. The explicit estimates
\[
\frac{e^{vt}-1}{t}\le ve^{vt},\qquad
\int_T^\infty\frac{e^{-\alpha t^2}}t\,dt
\le\frac{e^{-\alpha T^2}}{2\alpha T^2}
\] follow respectively by integrating the exponential derivative and by
using \(1/t\le t/T^2\). Choosing \(\alpha=\Delta_*/(2T)\) proves \[
\begin{split}
|\operatorname{Cov}_\psi(O_e,O_f)|
&\le\frac{2\|O_e\|\|O_f\|}{\pi}
        vT e^{-d(e,f)+vT}\\
&\quad+s_0\left(1+\frac2{\pi\Delta_*T}\right)
        e^{-\Delta_*T/2}.
\end{split}                                                   \tag{13.7}
\] The Gaussian is only an integral filter in this proof. The dynamics,
vacuum and interaction have never been replaced by Gaussian ones.

For \(d=d(e,f)\ge1\), now use \[
O_e=B_e^\Lambda,\quad O_f=B_f^\Lambda,\quad
\Lambda=e^{d/4},\qquad T=\frac{d}{\Delta_*+2v}.
\] The original parameters give exactly \[
\frac v{\Delta_*}=\frac{73728}{287}\xi\le\frac3{574},\qquad
vT\le\frac{3d}{580}<\frac d{100},\qquad
\Delta_*T\ge\frac{287d}{290}.                                \tag{13.8}
\] Consequently the first term in (13.7) is at most \[
\frac{73728}{287}\xi d\,e^{-49d/100}
\le1100\xi e^{-d/4}.
\] We used \(2/\pi<1\), and \(d e^{-6d/25}\le25/6\), which follows from
\(e^x\ge x\); the exact remaining coefficient is \(307200/287<1100\). By
(13.3), \(s_0\le64\xi^2\). The second term in (13.7) is at most
\(192\xi^2e^{-d/4}\): its parenthesis is less than \(3\) for \(d\ge1\)
and \(\Delta_*T/2\ge287d/580>d/4\). Finally (13.4) contributes at most
\(112\xi^2e^{-d/4}\). Their sum is bounded by
\((1100\xi+304\xi^2)e^{-d/4}\le1200\xi e^{-d/4}\). This proves (13.1)
for the original unbounded covariance.

\section{14. A volume-uniform two-sided remainder and the original
native
quotient}\label{a-volume-uniform-two-sided-remainder-and-the-original-native-quotient}

Define the real symmetric remainder \[
\mathscr R=C-\frac{\xi^2}{16}\mathsf I^{\mathsf T}\mathsf I.
\] For every entry (5.3) gives \(|\mathscr R_{ef}|\le42500\xi^3\). For
\(d(e,f)\ge2\) there is no shared face, so \(\mathscr R_{ef}=C_{ef}\)
exactly and (13.1) applies. These are two estimates of the same exact
entries; no correlation is set to zero.

The injective indexing map \(x(n,i)=2n+\mathbf e_i\) and its inverse
were proved in Section 9. Any two links on one face have indexing
\(\ell^1\) distance at most 2, as inspection of its four midpoints
shows. A graph ball of integer radius \(D\) therefore lies inside the
integer cube of coordinate side \(4D+1\), and contains at most
\((4D+1)^3\) links, including near the original box boundary. For a
shell of radius \(n\) we use the valid upper bound \((4n+1)^3\), not a
difference of two bounding cube volumes. For every integer \(D\ge1\),
this proves \[
\sup_e\sum_f|\mathscr R_{ef}|
\le42500\xi^3(4D+1)^3+
1200\xi\sum_{n>D}(4n+1)^3e^{-n/4}.                           \tag{14.1}
\]

Here is a closed, explicit bound with no remaining infinite sum. For
\(0<q<1\) put \[
G(A,q)=\frac{A^3}{1-q}
 +\frac{12A^2q}{(1-q)^2}
 +\frac{48Aq(1+q)}{(1-q)^3}
 +\frac{64q(1+4q+q^2)}{(1-q)^4}.
\] The geometric series and its first three applications of \(q\,d/dq\)
give respectively \[
\sum_{j\ge0}q^j=\frac1{1-q},\quad
\sum_{j\ge0}jq^j=\frac q{(1-q)^2},\quad
\sum_{j\ge0}j^2q^j=\frac{q(1+q)}{(1-q)^3},\quad
\sum_{j\ge0}j^3q^j=\frac{q(1+4q+q^2)}{(1-q)^4}.
\] Termwise differentiation is justified by uniform convergence on each
compact subinterval of \((0,1)\). Expanding \((A+4j)^3\) and using these
four identities proves \[
\sum_{n>D}(4n+1)^3q^n=q^{D+1}G(4D+5,q).                      \tag{14.2}
\] Write \(z=\log(1/\xi)>0\), set \(D=\lceil8z\rceil\), and use
\(q=e^{-1/4}<4/5\), the latter following from \(e^{1/4}>1+1/4\). Then
\(q^D\le\xi^2\), \(4D+1\le32z+5\), and \(4D+5\le32z+9\). All
coefficients of \(G\) are positive. Define the explicit cubic polynomial
\[
\mathcal P(z)=42500(32z+5)^3+1200G(32z+9,4/5),
\qquad \epsilon_\xi=\xi\mathcal P(\log(1/\xi)).              \tag{14.3}
\] Equations (14.1)--(14.2) now prove \[
\sup_e\sum_f|\mathscr R_{ef}|
\le\xi^3\mathcal P(\log(1/\xi))=\xi^2\epsilon_\xi.             \tag{14.4}
\] The same bound holds for column sums by symmetry. For completeness,
Cauchy--Schwarz in the nonnegative weights \(|\mathscr R_{ef}|\) gives
\[
\sum_e|(\mathscr Rw)_e|^2
\le\left(\sup_e\sum_f|\mathscr R_{ef}|\right)
 \sum_f|w_f|^2\sum_e|\mathscr R_{ef}|.
\] Thus (14.4) bounds its operator norm on \(\ell^2(\mathsf E_L)\) and
yields the two-sided inequality \[
\left|w^{\mathsf T}Cw-
          \frac{\xi^2}{16}\|\mathsf Iw\|_2^2\right|
\le\xi^2\epsilon_\xi\|w\|_2^2
\quad\text{for every real }w.                               \tag{14.5}
\] In particular the absolute rows of \(C\) are uniformly summable:
\(\sum_ft_{ef}=4r_e\le16\), so
\(\sup_e\sum_f|C_{ef}|\le\xi^2(1+\epsilon_\xi)\). The remainder has
order \(\xi^3(1+\log(1/\xi))^3\), uniformly in the full box. No
order-\(\xi^4\) assertion is substituted for that proved order.

There is an entirely explicit interval on which the relative estimate is
nonvacuous: \[
0<\xi\le10^{-16}\quad\Longrightarrow\quad\epsilon_\xi<1/32.
                                                               \tag{14.6}
\] To verify every constant, \(\mathcal P\) has positive coefficients
and degree three. For \(z\ge3\), every \(e^{-z}z^k\), \(0\le k\le3\), is
nonincreasing, by differentiation. Hence \(e^{-z}\mathcal P(z)\) is
nonincreasing there. The elementary bounds \(e<3<10<e^3\), proved by the
exponential series, give \(16<16\log10<48\). At \(z_0=16\log10\), \[
e^{-z_0}\mathcal P(z_0)
\le10^{-16}\mathcal P(48)
=\frac{356710369517}{20000000000000}<\frac1{32}.
\] The exact integer evaluation is \(\mathcal P(48)=178355184758500\).
Monotonicity proves (14.6), and the same exponential-series bounds show
\(\epsilon_\xi\to0\) as \(\xi\downarrow0\). The smaller interval is used
for the explicit positive lower denominators below; (14.5) holds on all
of (11.16).

For every nonzero nonnegative weight vector, summing the facewise
inequalities \((\sum w_e)^2\ge\sum w_e^2\) gives
\(\|\mathsf Iw\|_2^2\ge2\|w\|_2^2\). Equations (14.5)--(14.6) imply \[
\|v_w\|^2=w^{\mathsf T}Cw\ge\frac{3\xi^2}{32}\|w\|_2^2>0,
\qquad
\left|\frac{\mathfrak q_A[v_w]}{\|v_w\|^2}-3\kappa\right|
\le\kappa(3744000\xi+32\epsilon_\xi).                        \tag{14.7}
\] For the second inequality subtract \(3\kappa\) times (14.5) from
(7.2): the numerator error is at most
\(\kappa\xi^2(351000\xi+3\epsilon_\xi)\|w\|_2^2\). Divide by the
positive lower bound just proved. The conclusion is uniform over all
boxes and all such weights, not merely over a prescribed finite support.

For the unchanged native weights of (1.4), (14.5) reads \[
\left|\|v_{w^{\rm nat}}\|^2-\frac{\xi^2}{4}\mathcal A_L\right|
\le\xi^2\epsilon_\xi S_{2,L}
\le2\xi^2\epsilon_\xi\mathcal A_L.                            \tag{14.8}
\] Thus its relative denominator error is at most \(8\epsilon_\xi\),
with no \(L^3\) loss. On (14.6), \[
\frac{3\xi^2}{16}\mathcal A_L
\le\|v_{w^{\rm nat}}\|^2
\le\frac{5\xi^2}{16}\mathcal A_L,                             \tag{14.9}
\] and the exact physical native electric-state quotient satisfies
(14.7). No original \(n_2^2\) weight or original scale has been changed.
The physical substitution remains \(\xi=1/(4g_{\rm YM}^4)\),
\(\kappa=2g_{\rm YM}^2/a\); the explicit interval (14.6) is
\(g_{\rm YM}^4\ge2500000000000000\). This conservative numerical
interval is not claimed optimal.

For signed weights (14.5) is still valid, but the nonnegative-weight
lower bound is not asserted. On the full leading kernel it gives \[
0\le\|v_w\|^2\le\xi^2\epsilon_\xi\|w\|_2^2
\qquad(\mathsf Iw=0).                                       \tag{14.10}
\] The fourth-order forms already proved in the parent remain intact.
Equations (7.3) and (12.14) remain additional upper bounds, and the
smallest of these proved bounds may be used. A uniform fourth-order
kernel remainder and the minimum over all signed electric weights have
not been obtained by this argument. Section 15 below obtains a uniform
bound on that minimum by proving the gauge-equivariance of the full
cluster similarity and using the exact physical free threshold. It does
not identify the fourth-order kernel quotient by dividing (14.10).
Section 16 instead constructs the exact two-face correction function,
bounds its full-vacuum residual, and proves that kernel quotient with an
explicit uniform remainder in (16.24). Nor does (14.7) remove the
fixed-regulator restriction on the small-angle remainder in (1.4),
construct a spatial continuum theory, or determine the Millennium mass
gap.

\section{15. The gauge-restricted threshold and the full signed electric
minimum}\label{the-gauge-restricted-threshold-and-the-full-signed-electric-minimum}

The physical constraint improves the full-space estimate (12.8) by a
factor of four, without changing the interaction or the interval
(11.16). We prove the gauge maps, the free threshold, and the restricted
resolvent before using the improvement for the actual interacting
states.

The literature source for the finite-graph gauge Hilbert space and its
vertex-invariant representation is John C. Baez, \emph{Spin Network
States in Gauge Theory}, arXiv:gr-qc/9411007v1, the section \emph{Gauge
Theory on a Graph}, Lemmas 1, 2, and the lemma labelled lem2.5 in the
original TeX (\href{https://arxiv.org/abs/gr-qc/9411007v1}{primary
source}; \emph{Advances in Mathematics} 117 (1996), 253--272, DOI
10.1006/aima.1996.0012). The retained source is
\nolinkurl{literature/Baez-spin-network-states.tex}, lines 174--321. The
source constructs the invariant Hilbert space from edge representations
and vertex intertwiners. The leaf argument below proves the particular
free spectral bound directly, including the original open boundary and
every endpoint gauge transformation.

\subsection{15.1. Exact gauge projection, source convention, and support
sectors}\label{exact-gauge-projection-source-convention-and-support-sectors}

Retain the entire vertex set and define \[
\mathcal G_L=SU(2)^{\mathsf V_L},\qquad
(\alpha_gU)_e=g_{s(e)}U_eg_{t(e)}^{-1},\qquad
(\mathsf T_g f)(U)=f(\alpha_{g^{-1}}U).                       \tag{15.1}
\] There is no restriction of \(g_v\) to the identity at boundary
vertices and no external charged representation at a vertex. These are
precisely the gauge transformations of Section 1. Group multiplication
gives \(\alpha_g\alpha_h=\alpha_{gh}\), hence
\(\mathsf T_g\mathsf T_h=\mathsf T_{gh}\). Left and right Haar
invariance on each link prove \(\|\mathsf T_g f\|=\|f\|\) and
\(\mathsf T_g^*=\mathsf T_{g^{-1}}\). Translations of continuous
functions are continuous in \(g\), and density extends this to strong
continuity on \(\mathcal H_L\). Thus the strong Haar integral \[
\mathsf P_{\mathcal G}=\int_{\mathcal G_L}\mathsf T_g\,dg,
\qquad
\mathcal H_{\mathrm{phys}}
 =\{f:\mathsf T_gf=f\text{ for every }g\}
 =\operatorname{Ran}\mathsf P_{\mathcal G}                    \tag{15.2}
\] exists and is bounded. Haar invariance gives
\(\mathsf T_h\mathsf P_{\mathcal G}=\mathsf P_{\mathcal G}\); inversion
invariance gives \(\mathsf P_{\mathcal G}^*=\mathsf P_{\mathcal G}\).
Integrating the first identity proves
\(\mathsf P_{\mathcal G}^2=\mathsf P_{\mathcal G}\). Conversely the
integral fixes every invariant vector, proving both range identities in
(15.2).

Baez uses the coordinate action \(A_e\mapsto g_{t(e)}A_eg_{s(e)}^{-1}\).
The exact dictionary, on the same directed graph, is \[
\jmath:SU(2)^{\mathsf E_L}\longrightarrow SU(2)^{\mathsf E_L},
\qquad
A_e=\jmath(U)_e=U_e^{-1},\qquad \jmath^{-1}=\jmath.             \tag{15.3}
\] Indeed \((g_sU_eg_t^{-1})^{-1}=g_tU_e^{-1}g_s^{-1}\). The pullback
\(\mathsf Jf=f\circ\jmath\), from the source's Haar Hilbert space to
ours, is therefore a unitary intertwiner with inverse the same pullback.
Haar measure is invariant under inversion. Inversion is an isometry for
the bi-invariant metric with the \(T_a\) frame of (1.2), so it preserves
the Sobolev domains and intertwines every \(E_e\). This proves the
gauge-space dictionary; it does not silently reverse a Wilson traversal
in (1.1).

Each \(\mathsf T_g\) is a product of link isometries, so it preserves
\(H^1,H^2\) and commutes with \(E_e\) on its operator domain. In (1.1)
the endpoint factors cancel at the four consecutive vertices, leaving
conjugation by \(g_n\) inside the trace. Thus every \(W_p\), and
consequently \(H,H_0,V_p\), is gauge invariant on its stated domain. The
gauge projection also preserves these domains, by integrating in the
corresponding graph norm.

Let \(P_e\) be Haar averaging in the one link \(e\), followed by the
constant lift. Endpoint Haar invariance gives
\(P_e\mathsf T_g=\mathsf T_gP_e\). In Section 11 notation, the embedded
sector projection is exactly \[
\mathsf Q_I=
 \prod_{e\in I}(I-P_e)\prod_{f\notin I}P_f,
\qquad
\mathsf Q_If=(\Pi_If)\otimes1_{I^c}.                         \tag{15.4}
\] Products commute, so \(\mathsf Q_I\) commutes with the gauge
representation and with \(\mathsf P_{\mathcal G}\). The finite sum of
all \(\mathsf Q_I\) is the identity, by expanding
\(\prod_e[(I-P_e)+P_e]\). Restricting the unitary (11.2) consequently
gives the exact orthogonal decomposition \[
\mathcal H_{\mathrm{phys}}
 =\bigoplus_{I\subset\mathsf E_L}
     \bigl(\mathcal H'_I\otimes1_{I^c}\bigr)^{\mathcal G_L}.
                                                               \tag{15.5}
\] The inverse map remains the collection of Haar projections (15.4).
The action on the \(I\) factor is the restriction of the product link
representation, not a gauge group with its boundary removed.

\subsection{15.2. Leaf annihilation and the exact free physical
threshold}\label{leaf-annihilation-and-the-exact-free-physical-threshold}

For a nonempty set \(I\) of links, consider the undirected graph with
these edges and their endpoints. Suppose a vertex \(v\) has exactly one
incident edge \(e\) in \(I\). For \[
f\in\bigl(\mathcal H'_I\otimes1_{I^c}\bigr)^{\mathcal G_L},
\] vary only \(g_v\). On the \(I\) factors this translates only the
variable \(U_e\), on the left or the right according to its original
orientation. On the exterior factors it acts on constants. Since \(f\)
is invariant, integration over this single \(SU(2)\) gauge variable
gives \[
f=\int_{SU(2)}\mathsf T_{g_v}f\,dg_v=P_ef=0.                 \tag{15.6}
\] The second equality is exactly left or right Haar averaging; the last
equality follows from the \(I-P_e\) factor in (15.4). For arbitrary,
possibly entangled \(L^2\) vectors the identity follows first on
elementary tensors, then by density and the boundedness of the two
averages. Thus no pointwise representative or unstated smoothness is
used in the leaf argument.

It follows that every nonempty support of a nonzero invariant sector has
degree at least two at every one of its incident vertices. Such a finite
graph contains a cycle: choose a simple path of maximal length; at its
last vertex there is a neighbor other than the preceding vertex.
Maximality places that neighbor earlier in the path, and the resulting
segment closes a simple cycle. The original cubic graph has no loops or
multiple edges. The coloring \(n\mapsto(-1)^{n_1+n_2+n_3}\) changes sign
along every edge, so every cycle has even length. Its length is
therefore at least four. In particular, \[
\bigl(\mathcal H'_I\otimes1_{I^c}\bigr)^{\mathcal G_L}
 \ne\{0\},\quad I\ne\varnothing
\quad\Longrightarrow\quad |I|\ge4.                          \tag{15.7}
\] This proof includes vertices on faces, edges, and corners of the open
box. Omitting a boundary gauge transformation would invalidate the
corresponding step (15.6); no such omission has been made.

On a nonempty support (11.3) gives \(H_I\ge(3\kappa/4)|I|\), on its full
domain and hence on its invariant subspace. Combining (15.5)--(15.7) and
summing the orthogonal nonempty sectors yields the form inequality \[
H_0\big|_{\mathcal H_{\mathrm{phys}}\cap\Omega^\perp}
 \ge 3\kappa I.                                             \tag{15.8}
\] The constant is attained. For any elementary face \(p\), its
unchanged \(W_p\) is gauge invariant and is a fundamental matrix
coefficient in each of its four distinct link variables, including the
two inverse traversals. Bi-invariance of the Casimir gives \[
E_eW_p=\tfrac34W_p\quad(e\in\partial p),\qquad
E_eW_p=0\quad(e\notin\partial p).
\] Integration in one participating link, with all others held fixed,
gives \(\int W_p\,dU=0\) and \(\int |W_p|^2dU=1\). For the latter, the
product inside the trace is Haar distributed under that integral, also
for an inverse occurrence, and the fundamental matrix-coefficient
identity \(\int U_{ij}\overline{U_{k\ell}}\,dU
 =\delta_{ik}\delta_{j\ell}/2\) gives the unit norm of its trace.
Consequently \[
H_0W_p=3\kappa W_p,\qquad
W_p\in\mathcal H_{\mathrm{phys}}\cap\Omega^\perp,\qquad
\|W_p\|=1.                                                   \tag{15.9}
\] This proves the exact free physical threshold \(3\kappa\), rather
than merely a lower estimate. No free vector is substituted for the
interacting vacuum in the argument that follows.

\subsection{15.3. Gauge-equivariance of the interacting cluster
similarity}\label{gauge-equivariance-of-the-interacting-cluster-similarity}

The fixed-point construction of Section 11 may be performed within the
invariant sector collections. To prove that statement, assume each
\(z_I\otimes1_{I^c}\) is gauge invariant. The gauge operator factorizes
as \(\mathsf T_{g,I}\otimes\mathsf T_{g,I^c}\). Both \(z_I\) and \(1_I\)
are fixed by the first factor. Hence the rank-one creation operator
(11.4) satisfies \[
\mathsf T_g\widehat z_I\mathsf T_g^{-1}
 =|\mathsf T_{g,I}z_I\rangle
   \langle\mathsf T_{g,I}1_I|\otimes I_{I^c}
 =\widehat z_I.                                             \tag{15.10}
\] The projections \(\Pi_I\) intertwine the full representation with its
sector action by (15.4). The operator \(H_I^{-1}\) commutes with this
action by the reducing property and functional calculus. Since \(V_p\)
and \(\Omega\) are invariant, (11.10) therefore sends invariant
collections to invariant collections.

Iteration starts at zero. Every iterate is invariant, and the limit in
the energy norm (11.7) is invariant because fixed subspaces of bounded
representations are closed. Thus the actual fixed point has the claimed
invariance. Its creation sum \(C(z)\) and the two finite polynomial
exponentials commute with \(\mathsf T_g\) and
\(\mathsf P_{\mathcal G}\). The domain preservation proved in (11.9)
also holds upon restriction. In particular the exact relation remains \[
\psi=\alpha e^{C(z)}\Omega,\qquad
\alpha=\langle\Omega,\psi\rangle>0,
\] with the scalar and unit physical vacuum specified in (12.7).

Write \(\widetilde H=H_0+\mathcal K\) as in (12.1). All three operators
preserve the physical subspace; the bounded invertible similarity and
its inverse do as well. The restriction of the sector norm (12.2) is
still equivalent to the Hilbert norm, with its same explicitly
volume-dependent upper comparison. The relative estimate is unchanged:
\[
\|\mathcal K f\|_{\oplus,1}
 \le c_\xi\|H_0 f\|_{\oplus,1},\qquad
c_\xi=\frac{512}{3}\xi\le\frac1{288},
\quad
f\in\operatorname{Dom}(H_0)\cap\mathcal H_{\mathrm{phys}}.
                                                               \tag{15.11}
\] We do not replace the per-link bound used to prove this estimate by a
larger value or alter its numerical contraction.

On every nonempty invariant sector, the spectral parameter \(\lambda\)
of \(H_I\) satisfies \(\lambda\ge3\kappa\) by (15.7). Therefore, for
\(0<t<3\kappa(1-c_\xi)\), \[
\|\mathcal K(H_0-t)^{-1}\|_{\oplus,1;\mathrm{phys}}
 \le c_\xi\sup_{\lambda\ge3\kappa}
               \frac{\lambda}{\lambda-t}
 =\frac{3\kappa c_\xi}{3\kappa-t}<1.                         \tag{15.12}
\] The constant sector contributes zero because \(\mathcal K\Omega=0\),
exactly as in Section 12. The Neumann series therefore gives a bounded
inverse for \(\widetilde H-t\) on the physical subspace. Its image is in
the restricted \(H_0\) domain: factor the inverse as
\((H_0-t)^{-1}[I+\mathcal K(H_0-t)^{-1}]^{-1}\). Norm equivalence and
the invariant similarities transfer that inverse to \(H-\mathcal E-t\)
on the original physical Hilbert space. No form positivity is claimed
for the nonunitary similarity itself. Applying the spectral theorem to
the original self-adjoint restriction, with its simple ground
eigenvalue, proves \[
(H-\mathcal E)\big|_{\mathcal H_{\mathrm{phys}}\cap\psi^\perp}
 \ge 3\kappa\left(1-\frac{512}{3}\xi\right)I
 \ge\frac{287\kappa}{96}I,
\qquad 0<\xi\le\frac1{49152}.                                \tag{15.13}
\] In original parameters the last constant is \(287g_{\rm YM}^2/(48a)\)
and the range is \(g_{\rm YM}^4\ge12288\). The full, not necessarily
invariant Hilbert-space bound remains (12.8), with constant
\(287\kappa/384\). The fourfold improvement is asserted exactly on the
physical restriction just proved.

\subsection{15.4. Covariance order, signed states, and an attained
uniform
minimum}\label{covariance-order-signed-states-and-an-attained-uniform-minimum}

Every \(E_e\) commutes with the gauge representation, so each centered
vector \(v_e\) is in \(\mathcal H_{\mathrm{phys}}\cap\psi^\perp\). There
is no sign restriction on the coefficients in a real linear combination.
Applying (15.13) to every \(v_w\), including a zero vector, gives the
matrix inequality \[
\mathcal N\ \ge\
3\kappa(1-c_\xi)C\ \ge\ 0
\quad\hbox{on }\mathbb R^{\mathsf E_L}.                       \tag{15.14}
\] Combining it with the upper side of (7.2) yields \[
0\le w^{\mathsf T}Cw
 \le\frac{\xi^2\|\mathsf Iw\|_2^2/16
               +117000\xi^3\|w\|_2^2}{1-512\xi/3}.           \tag{15.15}
\] Here \(351000/3=117000\); the leading coefficient is \(1/16\). The
exact facewise estimate \(\|\mathsf Iw\|_2^2\le16\|w\|_2^2\) from
Section 12 gives \[
0\le C\le
\frac{\xi^2+117000\xi^3}{1-512\xi/3}
 I_{\mathbb R^{\mathsf E_L}}.                                \tag{15.16}
\] On \(\ker\mathsf I\) the numerator of (15.15) is exactly
\(117000\xi^3\|w\|_2^2\). This is an upper bound, not a fourth-order
identification or a discarded null direction. For the unchanged native
weights, \[
\|v_{w^{\rm nat}}\|^2
 \le\frac{\xi^2\mathcal A_L/4+117000\xi^3 S_{2,L}}
              {1-512\xi/3}
 \le\frac{\xi^2\mathcal A_L}{4}
       \frac{1+936000\xi}{1-512\xi/3}.                       \tag{15.17}
\] The last step uses \(S_{2,L}\le2\mathcal A_L\), without changing the
coefficients or coordinates in the original state.

To define the signed minimum without dividing by an unproved
denominator, let \(\mathcal H_{\mathrm{phys},\mathbb R}\) denote the
real-valued invariant functions and define \[
\mathsf V_\xi:\mathbb R^{\mathsf E_L}
 \longrightarrow\mathcal H_{\mathrm{phys},\mathbb R}\cap\psi^\perp,
\qquad
\mathsf V_\xi w=v_w,\qquad
\mathscr W_\xi=\operatorname{Ran}\mathsf V_\xi.               \tag{15.18}
\] Its Gram identity is
\(\langle\mathsf V_\xi w,\mathsf V_\xi z\rangle=w^{\mathsf T}Cz\). Thus
\(\ker\mathsf V_\xi=\ker C\): \(Cw=0\) implies zero norm of
\(\mathsf V_\xi w\); conversely a zero vector has zero inner product
with every \(\mathsf V_\xi z\). The induced map \[
\overline{\mathsf V}_\xi:
\bigl(\mathbb R^{\mathsf E_L}/\ker C,\,
       \langle[w],[z]\rangle_C=w^{\mathsf T}Cz\bigr)
 \longrightarrow\mathscr W_\xi                              \tag{15.19}
\] is an isometric linear bijection. The bilinear form is well-defined
because \(C\) annihilates its kernel, and it is positive on every
nonzero equivalence class by the Gram identity. The inverse sends a
vector \(v_w\) to its class \([w]\); two preimages differ by \(\ker C\),
which proves independence of choice. The same identity for the energy
form shows that \(\mathcal N\) annihilates \(\ker C\) and descends to
the quotient. This records the exact information loss. No equality
\(\ker C=\ker\mathsf I\), or \(\ker C=\{0\}\), is assumed.

On \(0<\xi\le1/680000\), (7.5)--(7.6) give a nonzero \(v_e\) for every
link, so \(\mathscr W_\xi\ne\{0\}\). It is finite dimensional and
consists of smooth vectors. The energy form is a continuous quadratic
form on this finite dimensional space. Its restriction to the unit
sphere attains a minimum by compactness. Through (15.19) this number is
exactly \[
\mathfrak m_L(\xi)
 =\min_{\substack{w\in\mathbb R^{\mathsf E_L}\\
                  w^{\mathsf T}Cw>0}}
       \frac{w^{\mathsf T}\mathcal N w}{w^{\mathsf T}Cw}.
                                                               \tag{15.20}
\] Write \(\Delta_{\mathrm{phys},L}(\xi)\) for the first positive
spectral value of \(H-\mathcal E\) on \(\mathcal H_{\mathrm{phys}}\).
Its existence follows from compact resolvent, the simple ground
eigenvalue, and the nonzero physical excited subspace; this is the same
excitation operator as in (15.13). The spectral variational principle,
(15.14), and the actual single-link trial bound (7.6) prove the complete
chain \[
\begin{split}
3\kappa-512\kappa\xi
 &=3\kappa(1-c_\xi)\\
 &\le\Delta_{\mathrm{phys},L}(\xi)
 \le\mathfrak m_L(\xi)
 \le3\kappa+2472000\kappa\xi ,
\qquad 0<\xi\le\frac1{680000}.
\end{split}                                                \tag{15.21}
\] Every \(L\ge2\) is covered by these same constants and coupling
interval. In particular the estimate permits weights that depend on
\(L,\xi\), have arbitrary signs, or lie in \(\ker\mathsf I\), whenever
they represent a nonzero actual state. Such a state is included in the
minimization, not excluded by a leading-coefficient test. The quotient
construction retains all zero actual states as the explicitly identified
kernel, where an energy quotient has no defined value.

Equation (15.21) proves both uniform statements
\(\Delta_{\mathrm{phys},L}(\xi)=3\kappa+O(\kappa\xi)\) and
\(\mathfrak m_L(\xi)=3\kappa+O(\kappa\xi)\), with the displayed
two-sided constants as their meaning. It does not assert that every
signed state's quotient is near \(3\kappa\). In particular it does not
replace the parent's fixed-box fourth-order kernel quotient or prove its
remainder uniform. The separate calculation (16.17)--(16.24) supplies
that uniform remainder on the entire original incidence kernel. The
earlier two-sided nonnegative/native estimate (14.7) also remains valid
with its original range. In physical parameters the range of (15.21) is
\(g_{\rm YM}^4\ge170000\) and its central scale is
\(3\kappa=6g_{\rm YM}^2/a\). No lattice-spacing limit or coupling
continuation is taken.

\subsection{15.5. The physical conditional-complement
operator}\label{the-physical-conditional-complement-operator}

Retain \(J_S,Q_S^\rho,\mathscr K_S^\rho,U_\psi,q_\rho\) and their
domains from Section 10. Let the gauge group act on \(\mathscr H_\rho\)
by the same pullback (15.1). It is unitary because \(\rho=\psi^2\) is
invariant. The retained-variable space has the induced endpoint action
on its retained links. Under a simultaneous transformation of \((s,r)\),
invariance of \(\rho(s,r)\) and Haar measure in \(s\) proves invariance
of the marginal \(\rho_{S^c}(r)\). Changing variables in the conditional
integral then gives \[
J_S^*\mathsf T_g=\mathsf T_{g,S^c}J_S^*,\qquad
\mathsf T_gJ_S=J_S\mathsf T_{g,S^c}.                          \tag{15.22}
\] It follows that \(Q_S^\rho\) and the gauge projection commute. They
preserve the form domain as proved in Section 10 and by the isometric
gauge action. Averaging smooth approximants over the compact gauge group
and then applying \(Q_S^\rho\) proves form density of the smooth
invariant conditional-mean-zero vectors in \[
\mathscr K_{S,\mathrm{phys}}^\rho
 =\mathscr K_S^\rho\cap\mathscr H_\rho^{\mathcal G_L}.
\] The closed form restricted to
\(H^1\cap\mathscr K_{S,\mathrm{phys}}^\rho\) therefore defines a
self-adjoint operator \(C_{S,\mathrm{phys}}^\rho\) by exactly the
variational domain rule given for \(C_S^\rho\) in Section 10. Equations
(15.22) and invariance of \(q_\rho\) also show that this is the reducing
physical restriction of \(C_S^\rho\).

The exact unitary \(U_\psi G=\psi G\) intertwines the gauge
representations. For a conditional-mean-zero \(G\), Fubini gives
\(\int G\rho\,dU=0\); consequently \(U_\psi G\) is physical and
perpendicular to \(\psi\) when \(G\) is invariant. Apply (15.13) under
this unitary and the form identity (10.10): \[
C_{S,\mathrm{phys}}^\rho
 \ge3\kappa(1-512\xi/3)I\ge\frac{287\kappa}{96}I,
\qquad
\|(C_{S,\mathrm{phys}}^\rho)^{-1}\|
 \le\frac1{3\kappa(1-512\xi/3)}
 \le\frac{96}{287\kappa}.                                    \tag{15.23}
\] This holds for every nonempty eliminated set \(S\), independently of
its size and the exterior volume, on (11.16). It concerns the full form
with derivatives in all link directions, not a differential operator
containing only the eliminated derivatives. For noninvariant
conditional-mean-zero vectors the previously proved constants
(12.15)--(12.16), not (15.23), remain the stated bounds. The
mathematical distinction is implemented by the proved reducing
projection and unitary, not by declaring the two constructions
unrelated.

\section{16. Uniform fourth-order forms on the entire signed incidence
kernel}\label{uniform-fourth-order-forms-on-the-entire-signed-incidence-kernel}

Retain the full Hamiltonian (1.2), its actual positive unit vacuum
\(\psi\), \(A=H-\mathcal E\), the gauge constraint (15.1), and every
original real edge weight. Throughout this section \(0<\xi\le1/49152\),
unless an explicitly smaller interval is displayed. Write
\(\mathscr Z_L=\ker\mathsf I\). We prove volume-uniform remainders for
the first surviving covariance and energy forms on all of
\(\mathscr Z_L\), not merely on one alternating weight.

The exact two-face coefficients are the existing calculation in the
parent source
\nolinkurl{wilson_variation/ELECTRIC_COVARIANCE_UNSIGNED_INCIDENCE_20260908.md},
Sections 4--5, equations (4.1)--(5.8). The source's electric operator
denoted \(H_0\) there is \(\sum_eE_e\); our \(H_0\) in (11.1) is
\(\kappa\sum_eE_e\). Accordingly define here
\(\mathfrak h_0=\sum_eE_e\), so \(H_0=\kappa\mathfrak h_0\), on the same
domain. Every occurrence of a physical energy below retains
\(\kappa=2g_{\rm YM}^2/a\). No change of Hamiltonian or vacuum is made.
We first record the two-face functions and their exact derivative
dictionary needed for the new remainder calculation.

\subsection{16.1. The retained two-face modes and a local correction
function}\label{the-retained-two-face-modes-and-a-local-correction-function}

Let \(\mathscr P_2\) be the unordered pairs of distinct elementary faces
sharing a link. Their unique shared link is \(e(p,q)\); their other six
links form the set \(O(p,q)\). Write \[
R_{pq}=\operatorname{tr}(A_{pq}B_{pq}),\qquad
D_{pq}=W_pW_q-\tfrac12R_{pq},\qquad
G_{pq}=\nabla_{e(p,q)}W_p\cdot\nabla_{e(p,q)}W_q .
\] Here the original face words, cyclically based at the shared link,
are \(U_eA_{pq}\) and \(U_e^{-1}B_{pq}\), reversing a full face
traversal when necessary. The trace is unchanged under that full
reversal in \(SU(2)\). The words \(A_{pq},B_{pq}\) contain the six
distinct outer links with the original inherited inverse occurrences. No
individual inverse occurrence is dropped.

The Pauli contraction with \(T_a=-i\sigma_a/2\) is \[
\sum_a\operatorname{tr}(T_aC)\operatorname{tr}(T_aD)
 =-\tfrac12\operatorname{tr}(CD)
   +\tfrac14\operatorname{tr}C\operatorname{tr}D .
\] Inserting the minus sign from differentiation of the inverse shared
link proves the first identity below. The product rule for \(E_e\) then
proves the shared-link label \(2\) of \(D_{pq}\): \[
\begin{split}
G_{pq}&=\tfrac38R_{pq}-\tfrac14D_{pq},\\
\mathfrak h_0R_{pq}&=\tfrac92R_{pq},\qquad
\mathfrak h_0D_{pq}=\tfrac{13}{2}D_{pq},\\
\|R_{pq}\|_0^2&=1,\qquad
\|D_{pq}\|_0^2=\tfrac34,\qquad
\langle R_{pq},D_{pq}\rangle_0=0 .
\end{split}                                                  \tag{16.1}
\] The subscript \(0\) denotes the original full product Haar measure,
used to compute coefficients, not substituted for the physical density.
Indeed each outer link has Casimir \(3/4\) on both functions; on the
shared link \(R_{pq}\) has Casimir zero. The identity
\(E_e(W_pW_q)=2D_{pq}\) follows from the displayed Pauli contraction,
and gives the remaining Casimir label.

For completeness, the norms in (16.1) can be checked by conditioning on
the shared link. The two disjoint outer words are independent Haar
matrices. Specifically take \(P=U_eA_{pq}\) and \(Q=B_{pq}U_e^{-1}\);
then \(\operatorname{tr}(PQ)=R_{pq}\) and
\(\operatorname{tr}P\operatorname{tr}Q=W_pW_q\). For these independent
holonomies \(P=p_0I+i p_a\sigma_a,\ Q=q_0I+i q_a\sigma_a\), put
\(x=p_0q_0,\ y=\sum_{a=1}^3p_aq_a\). Their uniform unit-sphere second
moments give
\(\mathbb E x^2=1/16,\ \mathbb E y^2=3/16,\ \mathbb E xy=0\). Now
\(R=2(x-y)\) and \(D=3x+y\); multiplying these expressions proves all
three Haar products in (16.1). Their individual Haar means are zero, by
integration in an outer fundamental link.

The functions \(R_{pq},D_{pq}\), for all unordered pairs, are mutually
orthogonal except for their own squared norms. Here is the geometric
check retained from the parent. A coplanar pair has a \(2\)-by-\(1\)
rectangular outer boundary, which determines its two unit faces. A
perpendicular pair has a unit-cube bounding box; the two vertices absent
from its outer six-cycle form the edge opposite the shared edge. They
determine that shared edge and the two faces. The two cases have
respectively two and three nonzero coordinate spans. Thus different
pairs have different outer supports. Two \(R\)'s then have a
\(3/4\)-versus-zero link label somewhere. An \(R\) and a \(D\) differ at
the latter's label-\(2\) link. Two \(D\)'s differ either at that
label-\(2\) link or at an outer link. Self-adjointness of the differing
link Casimir proves the claimed orthogonality in every case.

For a physical link \(e\) and pair \(\{p,q\}\), define the exact real
coefficient \[
\begin{split}
c_{e;pq}=-\mathbf1_{\{e=e(p,q)\}}
              +\tfrac16\mathbf1_{\{e\in O(p,q)\}},&\\
F_e=\sum_{\{p,q\}\in\mathscr P_2}c_{e;pq}G_{pq},\qquad&\\
\mathcal B_e=\sum_{\{p,q\}\in\mathscr P_2}c_{e;pq}
          \left(\tfrac1{12}R_{pq}-\tfrac1{26}D_{pq}\right).&
\end{split}
                                                               \tag{16.2}
\] Equation (16.1) proves \(\mathfrak h_0\mathcal B_e=F_e\) and
\(\int\mathcal B_e\,dU=0\). These functions are bounded, real, smooth
and gauge invariant. They are not truncated Hamiltonians.

Every contributing pair has \(e\) in at least one of its faces. There
are at most four such faces and at most twelve adjacent faces for each,
giving at most \(48\) pairs. Their link union lies in the shared-face
graph ball of radius two about \(e\). Each radius-one ball has at most
\(13\) links, so a radius-two ball is contained in a union of \(13\)
such balls and has at most \(169\) links. The pointwise bounds are \[
\|\mathcal B_e\|_\infty\le24,\qquad
\sum_f\|\nabla_f\mathcal B_e\|_\infty\le168.                  \tag{16.3}
\] To verify the constants, \(|R|\le2,\ |D|\le5\). On an outer link
\(|\nabla R|\le1\) and \(|\nabla D|\le2+1/2=5/2\), using (6.2) and the
product rule. On the shared link \(\nabla R=0,\ |\nabla D|\le4\). The
absolute value of a shared-coefficient summand in (16.2) is at most
\(2/12+5/26<1/2\). Its derivative at any one of its seven links is at
most \(\max\{1/12+(5/2)/26,4/26\}<1/2\). An outer-coefficient summand is
a further factor \(1/6\) smaller. Summing at most \(48\) summands and
seven derivative locations proves (16.3).

\subsection{16.2. A full-vacuum third-order state
residual}\label{a-full-vacuum-third-order-state-residual}

Set \(h_e=W_e^\star/4\), which is a function, and put
\(d_e=E_e\psi-\xi h_e\psi\). Section 4 gives \(\|d_e\|\le4200\xi^2\). We
need its differentiated counterpart, not a derivative of a norm
estimate. Introduce \(f_e=\psi-(\xi/3)W_e^\star\psi\). The exact product
rule gives \[
E_ef_e=d_e-\tfrac{\xi}{3}W_e^\star E_e\psi
                         +\tfrac{2\xi}{3}D_e\psi .
\] Here \(D_e\) is the differential contraction of Section 6, not the
two-face function \(D_{pq}\). Equations (3.1), (10.5) and (6.2) give
\(\|E_e\psi\|\le8\xi,\ \|\nabla_e\psi\|\le8\xi,\
\|D_e\psi\|\le32\xi\). Thus
\(\|E_ef_e\|\le(4200+128/3)\xi^2=(12728/3)\xi^2\). The exact one-link
spectral inequality \(\|\nabla_ef\|\le(2/\sqrt3)\|E_ef\|\) holds also
when \(f\) has a constant component. Apply it to \(f_e\) and retain the
product-rule remainder: \[
\begin{split}
\left\|\nabla_e\psi-\tfrac{\xi}{3}
                    (\nabla_eW_e^\star)\psi\right\|
&\le\frac{2}{\sqrt3}\frac{12728}{3}\xi^2
                              +\frac{64}{3}\xi^2\\
&<5000\xi^2 .
\end{split}                                                  \tag{16.4}
\] For the last numerical inequality use \(2/\sqrt3<7/6\); the resulting
coefficient is \(44740/9<5000\). All differentiations are on the actual
smooth vacuum.

From (6.1) and the full ground-state product rule, \[
Ad_e=\kappa\xi\left[
 -2D_e\psi+\tfrac12
   \sum_{p\ni e}\sum_{f\in\partial p}
          \nabla_fW_p\cdot\nabla_f\psi\right].
                                                               \tag{16.5}
\] Insert (16.4) for each differentiated vacuum. The error from the
first term has norm at most \(2\cdot4\cdot5000\kappa\xi^3\). The error
from the second has norm at most
\((1/2)\cdot4\cdot4\cdot5000\kappa\xi^3\). The full leading function is
exactly \(F_e\) of (16.2). To check every cancellation, a repeated face
contributes \(-2/3\) from the first sum and \(4/6\) from the second,
hence zero. An unordered distinct pair with shared link \(e\)
contributes \(-4/3+2/6=-1\); a pair with \(e\) on its outer boundary
contributes \(1/6\). These are precisely \(c_{e;pq}\). Consequently \[
\|Ad_e-\kappa\xi^2F_e\psi\|\le80000\kappa\xi^3.               \tag{16.6}
\]

Define the original local differential observable and its centered
residual state by \[
\mathscr D_e=E_e-\xi h_e-\xi^2\mathcal B_e,\qquad
s_e=Q_\psi\mathscr D_e\psi,\qquad
Q_\psi=I-|\psi\rangle\langle\psi|.                            \tag{16.7}
\] It is a bounded real multiplication perturbation of \(E_e\), with the
same one-link operator domain tensored with the exterior space. It is
self-adjoint on that domain; every product below is first evaluated on
smooth \(\psi\). The vector is gauge invariant, real, smooth and
perpendicular to the same \(\psi\). Since \[
A(\mathcal B_e\psi)
 =\kappa F_e\psi
       -2\kappa\sum_f\nabla_f\mathcal B_e\cdot\nabla_f\psi ,
\] the last term has norm at most
\(2\kappa\cdot168\cdot8\xi=2688\kappa\xi\). Equations (16.6)--(16.7),
including the exact centering killed by \(A\), imply \[
\|As_e\|\le82688\kappa\xi^3<83000\kappa\xi^3,\qquad
\|s_e\|\le28000\xi^3.                                       \tag{16.8}
\] The second assertion is the proved physical inverse bound
\(96/(287\kappa)\) from (15.13), applied to \(s_e\):
\(96\cdot83000/287<28000\). It is not obtained from an assumed locality
of the full vacuum.

\subsection{16.3. Exact finite-range residual energy and its extensive
consequence}\label{exact-finite-range-residual-energy-and-its-extensive-consequence}

The support of \(\mathscr D_e\) is inside the radius-two link ball about
\(e\). The commutator \([H,\mathscr D_f]\) has support inside the
radius-two ball about \(f\) as well. Indeed \([H,E_f]\) uses only the
faces at \(f\), while the commutator of \(H\) with either multiplication
function in (16.7) contains only that function's derivatives on its own
link support. All multiplication potentials commute with those
functions. No distant coefficient is generated by these exact
commutators.

Let \(\mathcal S_{ef}=\langle s_e,As_f\rangle\). Expanding the double
commutator on the smooth vacuum gives \[
\mathcal S_{ef}
 =\tfrac12\langle\psi,[\mathscr D_e,[H,\mathscr D_f]]\psi\rangle,
\qquad
\mathcal S_{ef}=0\quad\text{when }d(e,f)>4.                  \tag{16.9}
\] For the first identity no commutation of
\(\mathscr D_e,\mathscr D_f\) is assumed. The expansion equals
\(\langle\mathscr D_e\psi,A\mathscr D_f\psi\rangle+
 \langle\mathscr D_f\psi,A\mathscr D_e\psi\rangle\); the two terms are
equal because all functions and differential operators are real and
\(A\) is self-adjoint. Centering changes neither term. For the second
identity the two indicated supports are disjoint, so their differential
operators commute on the smooth domain.

The radius-four link ball has at most \(17^3=4913\) vertices by the
unchanged coordinate map of Section 14. Equation (16.8) gives
\(|\mathcal S_{ef}|\le28000\cdot83000\kappa\xi^6\). The symmetric
row-sum estimate used in (6.6) now gives, for every real weight \(w\),
with \(s_w=\sum_e w_es_e\), \[
0\le\mathfrak q_A[s_w]\le10^{14}\kappa\xi^6\|w\|_2^2,
\qquad
\|s_w\|\le6\cdot10^6\xi^3\|w\|_2.                           \tag{16.10}
\] The exact row coefficient is
\(4913\cdot28000\cdot83000=11417812000000<10^{14}\). Applying (15.13)
once more bounds the squared norm by \((96/287)10^{14}\xi^6\|w\|_2^2\);
\((96/287)10^{14}<(6\cdot10^6)^2\). This is the step which removes the
number of nonzero weights from the residual estimate. The physical
\(s_w\) and its full energy form remain unchanged.

\subsection{16.4. Uniform covariance and energy estimates for the
correction
functions}\label{uniform-covariance-and-energy-estimates-for-the-correction-functions}

Put \[
u_e=Q_\psi\mathcal B_e\psi,\qquad
\mathcal C^{B}_{ef}=\langle u_e,u_f\rangle,\qquad
\mathcal C^{B,0}_{ef}=\langle\mathcal B_e,\mathcal B_f\rangle_0 .
\] Each \(\mathcal B_e\) uses at most \(169\) links, and a product uses
at most \(338\). Equations (3.2)--(3.3), with \(32\sqrt{169}/3<139\) and
\(32\sqrt{338}/3<197\), give \[
|\langle\mathcal B_e\rangle_\psi|\le7000\xi,\qquad
|\mathcal C^B_{ef}-\mathcal C^{B,0}_{ef}|\le230000\xi .
                                                               \tag{16.11}
\] In detail, the first coefficient is at most
\(48\cdot139(1+139/49152)<7000\). For the second use the product bound
\(24^2=576\), and retain the product of the two means. The coefficient
is at most \[
2\cdot576\cdot197+
 \frac{2\cdot576\cdot197^2+7000^2}{49152}
 =\frac{175757179}{768}<230000.
\] This controls every pair locally; it does not yet sum all pairs.

We also require the large-distance bound \[
|\mathcal C^B_{ef}|\le10000e^{-d(e,f)/4}
\quad\text{when }d(e,f)\ge5.                                \tag{16.12}
\] Here is the support-dependent version of the method in Section 13.
For bounded \(O_X,O_Y\) on disjoint link sets, the exact recurrence
(13.6) has the initial term
\(2\|O_X\|\|O_Y\|\mathbf1_{X\cap Y\ne\varnothing}\). The first face in
an iterated chain has at most \(4|X|\) choices; each later face has at
most \(16\) choices. Reaching \(Y\) requires at least \(d(X,Y)\) faces.
The same strong evolution and factorial-remainder proof therefore gives
\[
\|[\alpha_t(O_X),O_Y]\|
 \le2|X|\|O_X\|\|O_Y\|
      e^{-d(X,Y)}(e^{192b|t|}-1).                            \tag{16.13}
\] The extra factor is the explicitly counted \(|X|\), not the volume of
the full box or an assumed single-link estimate for a multi-link
observable. This is the crossing-set construction of Nachtergaele--Sims,
arXiv:1410.8174v1, Section 3, equations fbd, iterandum and the complete
Step 4; the original proof has been read directly in the retained TeX.

For \(\mathcal B_e,\mathcal B_f\), take \(X,Y\) their actual supports
and set \(s=d(e,f)-4\ge1\), so \(d(X,Y)\ge s\). Repeat the explicitly
evaluated Gaussian filter (13.7), now with the \(|X|\) in (16.13), and
take \(T=s/(\Delta_*+2v)\), \(\Delta_*=287\kappa/384,\ v=192b\). The
spectral-integral proof is unchanged because its only observable
hypotheses are boundedness and self-adjointness. This is the method of
Nachtergaele--Sims, arXiv:math-ph/0506030v3, Section 3.2; its spectral
and Gaussian integral proof is read and retained, not inferred from a
title. The numerical estimates (13.8) give \(vT<s/100\) and
\(\Delta_*T\ge287s/290\). Since \(|X|\le169\) and
\(\|u_e\|\|u_f\|\le576\), the two terms of the filter are at most \[
\frac{576\cdot169}{100}s\,e^{-99s/100}
       +1728e^{-s/4}
 \le\left(\frac{576\cdot169}{74}+1728\right)e^{-s/4}.
\] We used \(s e^{-74s/100}\le100/74\), proved by \(\exp(x)\ge x\), and
\(2/\pi<1\). Finally \(e^{-s/4}=\exp(1)e^{-d(e,f)/4}<3e^{-d(e,f)/4}\);
the resulting coefficient \(337824/37<10000\) proves (16.12).

For \(d(e,f)>4\), the two correction functions have disjoint supports
and zero individual Haar means, so \(\mathcal C^{B,0}_{ef}=0\) by
Fubini. Use (16.11) up to integer distance \(D\) and (16.12) outside it.
With the generating function \(G\) already explicitly evaluated in
(14.2), define \[
\begin{split}
\mathcal P_B(z)&=230000(16z+5)^3+10000G(16z+9,4/5),\\
r_B(\xi)&=\xi\mathcal P_B(\log(1/\xi)).
\end{split}                                                  \tag{16.14}
\] Take \(D=\lceil4\log(1/\xi)\rceil\), which is at least four on the
present interval. The same coordinate shell count gives \[
\begin{split}
\sup_e\sum_f|\mathcal C^B_{ef}-\mathcal C^{B,0}_{ef}|
&\le230000\xi(4D+1)^3
       +10000\sum_{n>D}(4n+1)^3e^{-n/4}\\
&\le r_B(\xi).
\end{split}                                                  \tag{16.15}
\] Indeed \(e^{-D/4}\le\xi\), \(4D+1\le16\log(1/\xi)+5\), and
\(4D+5\le16\log(1/\xi)+9\); (14.2) evaluates the full tail. The
symmetric row and column bound also bounds the \(\ell^2\) operator norm,
by the weighted Cauchy--Schwarz proof following (14.4). Thus (16.15)
applies to arbitrary extensive real weights without an \(\ell^1\) loss.

For energy, retain the exact multiplication-state form \[
\mathcal E^B_{ef}
 =\mathfrak q_A[u_e,u_f]
 =\kappa\int \rho\sum_h
       \nabla_h\mathcal B_e\cdot\nabla_h\mathcal B_f\,dU .
\] Its Haar counterpart is \(\mathcal E^{B,0}_{ef}
 =\kappa\langle\mathcal B_e,\mathfrak h_0\mathcal B_f\rangle_0\), by
integration by parts. Both matrices vanish when \(d(e,f)>4\). The
multiplier in the integral has sup norm at most \(168^2=28224\) and
support size at most \(338\). Formula (3.3) bounds the coefficient of
its expectation difference by \[
2\cdot28224\cdot197(1+197/49152)
=\frac{1429097691}{128}<12000000.
\] The same radius-four row count proves, for all real \(w\), \[
|w^{\mathsf T}(\mathcal E^B-\mathcal E^{B,0})w|
\le6\cdot10^{10}\kappa\xi\|w\|_2^2 .                         \tag{16.16}
\] The exact row coefficient \(4913\cdot12000000=58956000000\) is
smaller than the stated constant.

\subsection{16.5. The full kernel identity and its uniform fourth-order
remainder}\label{the-full-kernel-identity-and-its-uniform-fourth-order-remainder}

For \(w\in\mathscr Z_L\), the two face sums vanish separately. Thus the
sum of its six outer weights at each adjacent pair is \(-2w_{e(p,q)}\).
Substituting in (16.2) gives the exact function identity \[
\mathcal B_w:=\sum_e w_e\mathcal B_e
 =\sum_{\{p,q\}\in\mathscr P_2}w_{e(p,q)}
         \left(-\tfrac19R_{pq}+\tfrac2{39}D_{pq}\right)
 =S_w .                                                     \tag{16.17}
\] Indeed the summed \(c_{e;pq}\) is
\(-w_{e(p,q)}+(1/6)(-2w_{e(p,q)})=-4w_{e(p,q)}/3\). This is precisely
the parent's surviving two-face function; the present calculation has
not chosen a different weight family or discarded one of its Casimir
components.

Define \[
\mathcal K_L(w,z)=\sum_e\binom{r_e}{2}w_ez_e,\qquad
d_0=\frac{196}{13689},\qquad e_0=\frac8{117}.
\] Since \(2\le r_e\le4\),
\(\|w\|_2^2\le\mathcal K_L(w,w)\le6\|w\|_2^2\). The orthogonality in
(16.1) gives exactly \[
\|\mathcal B_w\|_0^2=d_0\mathcal K_L(w,w),\qquad
\kappa\langle\mathcal B_w,\mathfrak h_0\mathcal B_w\rangle_0
=\kappa e_0\mathcal K_L(w,w).                                \tag{16.18}
\] The retained component sums are \(d_0=1/81+(3/4)(4/1521)\) and
\(e_0=(9/2)/81+(13/2)(3/4)(4/1521)\). Each pair is counted once through
its unique shared link.

The first-order term in (16.7) cancels exactly on \(\mathscr Z_L\),
because \(\sum_e w_eh_e=(1/4)\sum_p(\mathsf Iw)_pW_p=0\). Consequently
the actual centered state satisfies \[
v_w=\xi^2u_w+s_w,\qquad
u_w=Q_\psi\mathcal B_w\psi,\qquad
s_w=\sum_e w_es_e,\qquad w\in\mathscr Z_L.                   \tag{16.19}
\] Both terms on the right are physical vectors of the full interacting
Hilbert space, on the domains used above.

Set the explicit positive errors \[
\begin{split}
\varepsilon_C(\xi)
 &=r_B(\xi)+12\cdot10^6\xi\sqrt{6d_0+r_B(\xi)}
                              +36\cdot10^{12}\xi^2,\\
\varepsilon_N(\xi)
 &=6\cdot10^{10}\xi
   +2\cdot10^7\xi\sqrt{6e_0+6\cdot10^{10}\xi}
                              +10^{14}\xi^2.
\end{split}                                                  \tag{16.20}
\] Then, simultaneously for every original box and every real
\(w\in\mathscr Z_L\), \[
\begin{split}
\left|\|v_w\|^2-d_0\xi^4\mathcal K_L(w,w)\right|
 &\le\xi^4\varepsilon_C(\xi)\|w\|_2^2,\\
\left|\mathfrak q_A[v_w]
             -\kappa e_0\xi^4\mathcal K_L(w,w)\right|
 &\le\kappa\xi^4\varepsilon_N(\xi)\|w\|_2^2 .
\end{split}                                                  \tag{16.21}
\] Here is the complete error calculation. Equations (16.15),(16.18)
give \[
\left|\|u_w\|^2-d_0\mathcal K_L(w,w)\right|
\le r_B(\xi)\|w\|_2^2,\qquad
\|u_w\|\le\sqrt{6d_0+r_B(\xi)}\,\|w\|_2.
\] Expand the squared norm of (16.19). Its cross term is at most
\(2\xi^2\|u_w\|\|s_w\|\), and its residual square is \(\|s_w\|^2\).
Substitution of (16.10) yields exactly the first line of (16.20), with
no hidden weight-dependent constant.

For energy, (16.16),(16.18) give \[
\left|\mathfrak q_A[u_w]-\kappa e_0\mathcal K_L(w,w)\right|
 \le6\cdot10^{10}\kappa\xi\|w\|_2^2,\qquad
\mathfrak q_A[u_w]\le
 \kappa(6e_0+6\cdot10^{10}\xi)\|w\|_2^2.
\] The nonnegative form satisfies Cauchy--Schwarz: apply ordinary
Cauchy--Schwarz to \(A^{1/2}u_w,A^{1/2}s_w\). Thus the energy cross term
in (16.19) is at most
\(2\xi^2\sqrt{\mathfrak q_A[u_w]\mathfrak q_A[s_w]}\). Use (16.10),
whose square-root coefficient is \(10^7\). The residual energy
contributes \(10^{14}\kappa\xi^6\|w\|_2^2\). These are exactly the three
terms in \(\varepsilon_N\), proving the second line of (16.21).

Both errors tend to zero as \(\xi\downarrow0\). Indeed \(\mathcal P_B\)
is a fixed cubic with positive coefficients, so
\(\xi\mathcal P_B(\log(1/\xi))\to0\); each remaining displayed term then
tends to zero. More explicitly, (16.21) gives a covariance remainder
\(O(\xi^5(1+\log(1/\xi))^3)\|w\|_2^2\) and an energy remainder
\(O(\kappa\xi^5)\|w\|_2^2\), with constants independent of \(L,w\).
These real-coupling estimates do not assert that the fixed-box even
Taylor expansion has a nonzero fifth-order term. The parent's
\(O_L(\xi^6)\) expansion remains true on its stated fixed-box interval;
its constant is not silently promoted to a uniform one.

\subsection{16.6. Nonzero kernel states and the uniform two-face energy
quotient}\label{nonzero-kernel-states-and-the-uniform-two-face-energy-quotient}

The explicit interval \[
0<\xi\le10^{-24}                                             \tag{16.22}
\] has \(\varepsilon_C(\xi)<d_0/2\). To check this without a numerical
logarithm, put \(z=\log(1/\xi)\). The positive-coefficient cubic
\(\mathcal P_B\) has \(e^{-z}\mathcal P_B(z)\) nonincreasing for
\(z\ge3\); differentiate each monomial \(e^{-z}z^k\), \(0\le k\le3\). At
the largest allowed \(\xi\), \(24<24\log10<72\), using the
exponential-series bounds \(e<10<e^3\). The exact polynomial value is \[
\mathcal P_B(72)=437811569040000 .
\] It follows for all of (16.22) that \[
r_B(\xi)\le
\frac{437811569040000}{10^{24}}<10^{-9},\qquad
6d_0+r_B(\xi)<1.
\] The second and third terms of \(\varepsilon_C\) are therefore bounded
by \(12\cdot10^6/10^{24}\) and \(36\cdot10^{12}/10^{48}\). Their sum
with the displayed rational bound is less than \(10^{-9}<d_0/2\); every
number in this comparison is rational. This conservative interval is
explicit, not asserted optimal.

For every nonzero real \(w\in\mathscr Z_L\), (16.21) now proves \[
\|v_w\|^2\ge\xi^4(d_0-\varepsilon_C(\xi))\|w\|_2^2
 \ge\tfrac12d_0\xi^4\|w\|_2^2>0.                             \tag{16.23}
\] The exact state map (15.18) is therefore injective on the entire
incidence kernel throughout the same volume-independent interval. Its
restriction is a bijection onto its image, and (16.23) bounds its
inverse, in the original edge norm and physical state norm, by
\(\xi^{-2}\sqrt{2/d_0}\). Neither the weights nor the state has been
rescaled in that assertion.

Finally \(e_0/d_0=234/49\). Subtract \((234/49)\kappa\) times the first
line of (16.21) from the second before dividing by the strictly positive
bound (16.23). The leading \(\mathcal K_L(w,w)\) terms cancel with their
exact coefficients. The result is \[
\left|
 \frac{\mathfrak q_A[v_w]}{\|v_w\|^2}
          -\frac{234}{49}\kappa
\right|
\le
\kappa\,
\frac{\varepsilon_N(\xi)+(234/49)\varepsilon_C(\xi)}
     {d_0-\varepsilon_C(\xi)}
\quad
\left(0\ne w\in\mathscr Z_L,\ 0<\xi\le10^{-24}\right).
                                                               \tag{16.24}
\] Its right-hand coefficient tends to zero independently of the box and
of every signed, volume-dependent or coupling-dependent kernel weight.
This proves the uniform first surviving two-face quotient of the actual
states, retaining both free Casimir components only as coefficients of
the full-vacuum calculation.

The minimum over all signed actual states in (15.21) and the kernel
quotients in (16.24) describe the same state map on different explicitly
specified domains. Their central energies are \(3\kappa\) and
\((234/49)\kappa\); the difference \((234/49-3)\kappa=(87/49)\kappa>0\)
is explicit. No implication of unrelatedness is drawn: (16.17)--(16.19)
prove the exact cancellation map and the surviving physical vectors. The
present injectivity assertion is on \(\mathscr Z_L\); it is not a proof
that \(\ker C=\{0\}\) on the entire edge space for a uniform coupling
interval. Section 17 proves that stronger statement by controlling the
mixed one-face and two-face coefficient covariance; it does not infer it
from the kernel restriction. The physical parameters remain
\(\xi=1/(4g_{\rm YM}^4)\), \(\kappa=2g_{\rm YM}^2/a\). Interval (16.22)
is \(g_{\rm YM}^4\ge250000000000000000000000\). No spatial continuum
construction, interchange of limits, or Millennium conclusion is
asserted by these finite-regulator strong-coupling estimates.

\section{17. Mixed face coordinates and uniform injectivity on the full
edge
space}\label{mixed-face-coordinates-and-uniform-injectivity-on-the-full-edge-space}

We retain the same actual vacuum and all real edge weights, without
restricting to \(\ker\mathsf I\). The local residual construction of
Section 16 already applies to these weights. Its remaining leading
function contains both the one-face and two-face coordinates. We now
control their complete mixed covariance. Throughout \(0<\xi\le1/49152\);
the smaller interval used for a positive lower bound will be stated
separately.

\subsection{17.1. Exact coefficient map, image and
inverse}\label{exact-coefficient-map-image-and-inverse}

For each unordered adjacent pair \(\alpha=\{p,q\}\), put \[
f_p=W_p,\qquad
f_\alpha=\tfrac1{12}R_{pq}-\tfrac1{26}D_{pq},\qquad
b_0=\tfrac1{144}+\tfrac3{4\cdot676}=\frac{49}{6084}.
\] The index set is the disjoint union
\(\mathscr J_L=\mathsf P_L\sqcup\mathscr P_2\). These are real,
gauge-invariant smooth functions with zero Haar means, supported on at
most seven original links and of absolute value at most two. The pair
functions in fact have absolute value less than \(1/2\), by the
constants in (16.3). For different faces the Wilson functions are
orthogonal by integration in a link appearing in only one face. Each has
squared norm one. Every one-face function is orthogonal to every
\(R_{pq}\): their four- and six-link fundamental supports differ. It is
orthogonal to every \(D_{pq}\), whose shared link has Casimir two rather
than zero or \(3/4\). Together with (16.1), these facts prove the entire
Haar Gram matrix, including its mixed blocks: \[
\langle f_j,f_k\rangle_0=
\begin{cases}
1,&j=k\in\mathsf P_L,\\
b_0,&j=k\in\mathscr P_2,\\
0,&j\ne k.
\end{cases}                                                   \tag{17.1}
\]

For the original weight \(w\in\mathbb R^{\mathsf E_L}\), define \[
\begin{split}
z_p&=(\mathsf Iw)_p,\\
c_\alpha&=\tfrac16(z_p+z_q)-\tfrac43w_{e(p,q)},\\
\mathsf T_\xi w&=a,\qquad
a_p=\tfrac{\xi}{4}z_p,\quad a_\alpha=\xi^2c_\alpha,\\
F_w&=\sum_{j\in\mathscr J_L}a_j f_j
     =\xi h_w+\xi^2\mathcal B_w .
\end{split}                                                   \tag{17.2}
\] The formula for \(c_\alpha\) follows by summing the coefficients in
(16.2): the six outer weights sum to \(z_p+z_q-2w_e\). Thus no term of
\(\mathcal B_w\) is lost. The exact state identity for all weights is \[
v_w=Q_\psi F_w\psi+s_w .                                     \tag{17.3}
\]

Every link has at least two incident faces. Fix the lexicographically
first pair \(\alpha_e=\{p_e,q_e\}\) of them, using the original integer
face labels. On the whole coefficient space define \[
\begin{split}
\mathsf L_\xi:\mathbb R^{\mathscr J_L}&\longrightarrow
                         \mathbb R^{\mathsf E_L},\\
(\mathsf L_\xi a)_e
 &=\frac{a_{p_e}+a_{q_e}}{2\xi}
                   -\frac{3a_{\alpha_e}}{4\xi^2}.
\end{split}                                                   \tag{17.4}
\] Substitution of (17.2) gives
\((z_{p_e}+z_{q_e})/8-3c_{\alpha_e}/4=w_e\). Consequently
\(\mathsf L_\xi\mathsf T_\xi=I\) exactly, \(\ker\mathsf T_\xi=\{0\}\),
and \(\mathsf T_\xi\mathsf L_\xi\) is an idempotent with image
\(\operatorname{Ran}\mathsf T_\xi\). A coefficient array \(a\) lies in
this image exactly when \(a=\mathsf T_\xi\mathsf L_\xi a\). This is an
explicit system of linear equations, with inverse (17.4) on that image.
The powers of \(\xi\) belong to this recorded coefficient map; the
original physical state and its norm have not been changed.

Let \((\mathsf Jz)_\alpha=z_p+z_q\). Each face has at most twelve
adjacent faces, since its four links each meet at most three other
faces. Therefore \[
\|\mathsf Jz\|_2^2
\le2\sum_{\{p,q\}\in\mathscr P_2}(z_p^2+z_q^2)
\le24\|z\|_2^2 .
\] Using \(w_{e(p,q)}=(\mathsf Jz)_\alpha/8-3c_\alpha/4\) at every pair,
not only the selected pair, gives \[
\|w\|_2^2\le\mathcal K_L(w,w)
\le\tfrac34\|z\|_2^2+\tfrac98\|c\|_2^2.                      \tag{17.5}
\] Indeed \(|x+y|^2\le2|x|^2+2|y|^2\) gives the coefficients
\(24/32=3/4\) and \(18/16=9/8\). Write \(d_0=196/13689\) as before. The
exact Haar value is \[
\begin{split}
\mathcal H_\xi(w):=\|F_w\|_0^2
 &=\tfrac{\xi^2}{16}\|\mathsf Iw\|_2^2
                         +b_0\xi^4\|c\|_2^2,\\
\mathcal H_\xi(w)&\ge\tfrac{d_0}{2}\xi^4\|w\|_2^2 .
\end{split}                                                   \tag{17.6}
\] For the second inequality multiply (17.5) by \((d_0/2)\xi^4\). The
pair coefficients agree because \(9d_0/16=b_0\), and the face
coefficient is bounded by \(\xi^2/16\), since \(\xi\le1\) and
\(3d_0/8<1/16\). No inverse singular value of \(\mathsf I\) was used.

\subsection{17.2. The full actual-vacuum Gram
matrix}\label{the-full-actual-vacuum-gram-matrix}

Anchor each face \((n;i,j)\), \(i<j\), at link \((n,i)\); anchor each
pair at its unique shared link. At most two face indices and at most six
pair indices have the same anchor. All supports lie in the radius-one
link ball of their anchor. Hence the number of indices whose anchors lie
within distance \(D\) of a fixed anchor is at most \(8(4D+1)^3\).

Put \[
\mathsf G^\psi_{jk}
=\langle Q_\psi f_j\psi,Q_\psi f_k\psi\rangle,\qquad
\mathsf G^0_{jk}=\langle f_j,f_k\rangle_0 .
\] Each support has at most seven links, and two supports together have
at most fourteen. In (3.2)--(3.3) use \((32/3)\sqrt7<29\),
\((32/3)\sqrt{14}<40\). The means satisfy \[
|\langle f_j\rangle_\psi|
\le116\xi(1+29\xi)<117\xi.
\] The product has absolute value at most four. Keeping the product of
the means, the entrywise difference is bounded by \[
|\mathsf G^\psi_{jk}-\mathsf G^0_{jk}|
\le\bigl[320+(12800+13689)/49152\bigr]\xi<400\xi.              \tag{17.7}
\]

For anchor distance \(d\ge3\), the supports are disjoint and their
distance is at least \(s=d-2\ge1\). Apply the full-dynamics estimate
(16.13) with \(|X|\le7\) and \(\|f_j\|\|f_k\|\le4\), and the exact
Gaussian integral proof (13.7), retaining \(T=s/(\Delta_*+2v)\). Here
\(s_0\le4\). The two resulting terms are bounded by \[
\frac{28}{100}s e^{-99s/100}+12e^{-s/4}
\le\left(\frac{28}{74}+12\right)e^{-s/4}.
\] Since \(e^{1/2}<2\), this proves \[
|\mathsf G^\psi_{jk}|\le32e^{-d/4}\qquad(d\ge3).             \tag{17.8}
\] The constants follow from the same strong on-site evolution and
spectral integrals as in Section16.4, using the original
Nachtergaele--Sims source proofs cited there. No infinite-volume
Hamiltonian or tensorized correlated vacuum is assumed.

For \(d\ge3\), \(\mathsf G^0_{jk}=0\) also follows from disjoint
supports and zero Haar means. Define, with \(G\) from (14.2), \[
\begin{split}
\mathcal P_{\mathrm{mix}}(z)
 &=8\left[400(16z+5)^3+32G(16z+9,4/5)\right],\\
\rho_{\mathrm{mix}}(\xi)
 &=\xi\mathcal P_{\mathrm{mix}}(\log(1/\xi)).
\end{split}                                                   \tag{17.9}
\] Taking \(D=\lceil4\log(1/\xi)\rceil\ge2\), summing (17.7) inside and
(17.8) outside, with the anchor multiplicity eight, gives \[
\sup_j\sum_k|\mathsf G^\psi_{jk}-\mathsf G^0_{jk}|
\le\rho_{\mathrm{mix}}(\xi),\qquad
\|\mathsf G^\psi-\mathsf G^0\|_{\ell^2\to\ell^2}
\le\rho_{\mathrm{mix}}(\xi).                                \tag{17.10}
\] The summation is exactly
\(8[400\xi(4D+1)^3+32\sum_{n>D}(4n+1)^3e^{-n/4}]\). Equation (14.2)
evaluates its complete tail and \(e^{-D/4}\le\xi\) proves (17.10). The
symmetric row/column argument of (14.4) proves the operator-norm
statement. Thus the mixed blocks have been bounded, not set to zero in
the actual density.

For \(a=\mathsf T_\xi w\), put \[
\mathcal A_\xi(w)=\|a\|_2^2
 =\tfrac{\xi^2}{16}\|z\|_2^2+\xi^4\|c\|_2^2 .
\] Equations (17.1) and (17.10) imply \[
\left|\|Q_\psi F_w\psi\|^2-\mathcal H_\xi(w)\right|
\le\rho_{\mathrm{mix}}(\xi)\mathcal A_\xi(w),\qquad
b_0\mathcal A_\xi(w)\le\mathcal H_\xi(w).                    \tag{17.11}
\] In particular the following complete all-weight covariance error
keeps both coefficient arrays: \[
\begin{split}
\left|\|v_w\|^2-\mathcal H_\xi(w)\right|
&\le \rho_{\mathrm{mix}}(\xi)\mathcal A_\xi(w)\\
&\quad+12\cdot10^6\xi^3\|w\|_2
 \sqrt{\mathcal H_\xi(w)+\rho_{\mathrm{mix}}(\xi)\mathcal A_\xi(w)}\\
&\quad+36\cdot10^{12}\xi^6\|w\|_2^2 .
\end{split}                                                   \tag{17.12}
\] This is the squared norm expansion of (17.3), the exact bound
(17.11), and the original residual norm (16.10). It applies to signed,
volume-dependent and coupling-dependent weights, including mixtures
approaching \(\ker\mathsf I\).

\subsection{17.3. A positive uniform interval for every raw
weight}\label{a-positive-uniform-interval-for-every-raw-weight}

On \(0<\xi\le10^{-24}\), the positive-coefficient cubic in (17.9) gives
\(\rho_{\mathrm{mix}}(\xi)\le10^{-24}\mathcal P_{\mathrm{mix}}(72)\), by
the monomial monotonicity and \(24<24\log10<72\) used in Section16.6.
Direct rational evaluation of (17.9) gives
\(\mathcal P_{\mathrm{mix}}(72)=7044756359040<10^{13}\), so \[
\rho_{\mathrm{mix}}(\xi)<10^{-11}<b_0/4,\qquad
(24\cdot10^6\xi)^2<d_0/2 .                                 \tag{17.13}
\] Both numerical comparisons are rational and hold on the whole stated
interval. Equations (17.11), (17.13) show
\(\|Q_\psi F_w\psi\|\ge(3/4)\sqrt{\mathcal H_\xi(w)}\), since
\(1-\rho_{\mathrm{mix}}/b_0>3/4>(3/4)^2\). The second inequality in
(17.13) and (16.10) give \(\|s_w\|\le(1/4)\sqrt{d_0/2}\xi^2\|w\|_2\).
Use (17.6) in the reverse triangle inequality for (17.3). For every real
\(w\) the result is \[
\boxed{\quad
\|v_w\|^2\ge\frac{d_0}{8}\xi^4\|w\|_2^2,\qquad
C(\xi)\ge\frac{d_0}{8}\xi^4 I_{\mathbb R^{\mathsf E_L}}
\quad(0<\xi\le10^{-24}).\quad}                               \tag{17.14}
\] This is on the entire original edge space, uniformly in \(L\). It
strengthens, rather than substitutes for, the sharper kernel lower bound
(16.23).

The physical state map \(V_\xi:w\mapsto v_w\) of (15.18) therefore has
zero kernel on this interval and is bijective onto its actual image,
with inverse norm at most \(\xi^{-2}\sqrt{8/d_0}\). An explicit inverse
on that image is \[
w=C(\xi)^{-1}
      \bigl(\langle v_e,v_w\rangle\bigr)_{e\in\mathsf E_L}.
                                                               \tag{17.15}
\] Indeed the vector on the right before applying \(C^{-1}\) is
precisely \(Cw\). Inequality (17.14) proves existence of this
finite-matrix inverse, and the physical inverse-norm bound follows
directly by taking square roots in (17.14). No identification of these
vectors as individual eigenvectors is made. The original signed minimum
(15.21), kernel quotient (16.24), and full mixed covariance (17.12) all
refer to this same injective state map on their stated intervals. The
original spacing \(a>0\), coupling \(g_{\rm YM}\) and every Hamiltonian
term remain unchanged. This calculation does not take a spatial
continuum limit or continue beyond the displayed strong-coupling
interval.

\section{18. Complete mixed energy and the original electric-state Ritz
spectrum}\label{complete-mixed-energy-and-the-original-electric-state-ritz-spectrum}

Write \(A=H-\mathcal E\), with exactly the Hamiltonian, domains and
actual vacuum of Section 1. The coefficients \(z,c,a=\mathsf T_\xi w\),
the function \(F_w\), and the positive quadratic form
\(\mathcal H_\xi(w)\) are those of Section 17. We first work on
\(0<\xi\le1/49152\), and then give an explicit smaller interval for all
quotients and ordered Ritz values. All weights remain the original real
link arrays; none is replaced by a unit vector or a different choice of
physical coefficients.

\subsection{18.1. The complete coefficient energy
matrix}\label{the-complete-coefficient-energy-matrix}

For \(j,k\in\mathscr J_L\) define \[
\mathsf K^\psi_{jk}
 =\kappa\int\psi^2\sum_e\nabla_e f_j\cdot\nabla_e f_k\,dU,
\qquad
\mathsf K^0_{jk}
 =\kappa\int\sum_e\nabla_e f_j\cdot\nabla_e f_k\,dU .       \tag{18.1}
\] Equation (1.5), polarized over real smooth functions, proves that
\(\mathsf K^\psi\) is exactly the energy Gram matrix of
\(Q_\psi f_j\psi\). Centering contributes no derivative. Every function
and state here is smooth on the full compact link manifold, so both the
form and operator domains are available.

Haar integration by parts gives
\(\mathsf K^0_{jk}=\kappa\langle f_j,\mathfrak h_0f_k\rangle_0\). Each
\(W_p\) has free energy \(3\). For a pair, retaining the two distinct
eigenvalues in (16.1) gives \[
\langle f_\alpha,\mathfrak h_0 f_\alpha\rangle_0
 =\frac{9/2}{144}+\frac{(13/2)(3/4)}{676}=\frac1{26}.
\] The link-Casimir orthogonality proved in Sections16.1 and17.1 annuls
every other Haar entry, including the face/pair entries. Thus the entire
coefficient matrix and its original-weight form are \[
\begin{split}
\mathsf K^0_{jk}&=
\begin{cases}
3\kappa,&j=k\in\mathsf P_L,\\
\kappa/26,&j=k\in\mathscr P_2,\\
0,&j\ne k,
\end{cases}\\
\mathcal J_\xi(w)&:=\kappa^{-1}a^t\mathsf K^0a
 =\frac{3\xi^2}{16}\|z\|_2^2+\frac{\xi^4}{26}\|c\|_2^2,\\
3\mathcal H_\xi(w)&\le\mathcal J_\xi(w)
 \le q_*\mathcal H_\xi(w),\qquad q_*:=\frac{234}{49}.
\end{split}                                                   \tag{18.2}
\] Indeed \((1/26)/b_0=234/49\), with the unchanged \(b_0=49/6084\).
These are Haar coefficient calculations, not a replacement of the actual
matrix \(\mathsf K^\psi\).

For a face, \(\sum_e\|\nabla_e f_p\|_\infty\le4\), by (6.2). For a pair,
the seven derivative bounds preceding (16.3) each are less than \(1/2\);
their sum is less than \(7/2<4\). Consequently the real multiplier
\(M_{jk}=\sum_e\nabla_e f_j\cdot\nabla_e f_k\) has
\(\|M_{jk}\|_\infty\le16\). Its support has at most fourteen links.
Apply the actual local marginal estimate (3.3) with \(q_S<40\xi\). It
yields \[
|\mathsf K^\psi_{jk}-\mathsf K^0_{jk}|
 \le1280\kappa\xi(1+40\xi)<1300\kappa\xi .                 \tag{18.3}
\] If the anchors of \(j,k\) have link-graph distance greater than two,
their link supports are disjoint. Each summand in \(M_{jk}\) then
vanishes pointwise. Both energy entries, not merely their difference,
are zero. There are at most \(8(4\cdot2+1)^3=8\cdot9^3\) indices at
anchor distance at most two from any one anchor. The exact row
coefficient is \(1300\cdot8\cdot9^3=7581600<8000000\). Put \[
\chi_E(\xi)=8000000\xi.
\] The symmetric row/column estimate therefore proves \[
\begin{split}
\sup_j\sum_k|\mathsf K^\psi_{jk}-\mathsf K^0_{jk}|
 &\le\kappa\chi_E(\xi),\\
\|\mathsf K^\psi-\mathsf K^0\|_{\ell^2\to\ell^2}
 &\le\kappa\chi_E(\xi),\\
\left|\mathfrak q_A[Q_\psi F_w\psi]-\kappa\mathcal J_\xi(w)\right|
 &\le\kappa\chi_E(\xi)\mathcal A_\xi(w).
\end{split}                                                   \tag{18.4}
\] All actual mixed entries are retained in the left sides of
(18.3)--(18.4); their smallness was calculated from the vacuum.

Now use the exact identity (17.3), not an identification with its
leading part. Positivity of \(A\) and Cauchy--Schwarz for its form give
the complete energy error \[
\begin{split}
|w^t\mathcal Nw-\kappa\mathcal J_\xi(w)|
&\le\kappa\chi_E(\xi)\mathcal A_\xi(w)\\
&\quad+2\cdot10^7\kappa\xi^3\|w\|_2
 \sqrt{\mathcal J_\xi(w)+\chi_E(\xi)\mathcal A_\xi(w)}\\
&\quad+10^{14}\kappa\xi^6\|w\|_2^2.
\end{split}                                                   \tag{18.5}
\] The cross term is bounded by
\(2\sqrt{\mathfrak q_A[Q_\psi F_w\psi]\mathfrak q_A[s_w]}\); (18.4)
bounds its first factor and (16.10) its second. This proves (18.5) for
every weight, including any dependence of the weight on the box and
coupling.

\subsection{18.2. Uniform relative errors and every electric-state
quotient}\label{uniform-relative-errors-and-every-electric-state-quotient}

Keep \(d_0=196/13689\) and \(\rho_{\rm mix}\) from (17.9). Define the
following explicit functions of the original \(\xi\): \[
\begin{split}
r(\xi)&=\rho_{\rm mix}(\xi)/b_0,\qquad
t_C(\xi)=6\cdot10^6\sqrt{2/d_0}\,\xi,\\
t_N(\xi)&=10^7\sqrt{2/d_0}\,\xi,\\
\beta_C(\xi)&=r(\xi)+2t_C(\xi)\sqrt{1+r(\xi)}+t_C(\xi)^2,\\
\beta_N(\xi)&=\frac{\chi_E(\xi)}{b_0}
 +2t_N(\xi)\sqrt{q_*+\frac{\chi_E(\xi)}{b_0}}+t_N(\xi)^2.
\end{split}                                                   \tag{18.6}
\] By (17.6) and (17.11),
\(\|w\|_2\le\sqrt{2/d_0}\xi^{-2}\sqrt{\mathcal H_\xi(w)}\) and
\(\mathcal A_\xi(w)\le\mathcal H_\xi(w)/b_0\). Substitution into every
term of (17.12) and (18.5), followed by the two inequalities in (18.2),
proves \[
\begin{split}
|w^tCw-\mathcal H_\xi(w)|
 &\le\beta_C(\xi)\mathcal H_\xi(w),\\
|w^t\mathcal Nw-\kappa\mathcal J_\xi(w)|
 &\le\kappa\beta_N(\xi)\mathcal H_\xi(w).
\end{split}                                                   \tag{18.7}
\] For example the covariance cross term, divided by
\(\mathcal H_\xi(w)\), is at most
\(12\cdot10^6\sqrt{2/d_0}\xi\sqrt{1+r}\), exactly the middle term of
\(\beta_C\). The energy calculation has the same retained two factors
and yields its displayed \(\beta_N\). All these functions tend to zero
as \(\xi\downarrow0\): \(\rho_{\rm mix}\) is \(\xi\) times a fixed cubic
in \(\log(1/\xi)\), and all other factors are displayed.

On the explicit nonempty interval \(0<\xi\le10^{-24}\), set
\(x_0=10^{-24}\) only for recording numerical bounds. Equation (17.13),
\(2/d_0=13689/98<144\), and rational comparisons give \[
\begin{split}
\beta_C(\xi)&<B_C:={10^{-11}\over b_0}
 +4(72\cdot10^6x_0)+(72\cdot10^6x_0)^2<10^{-8},\\
\beta_N(\xi)&<B_N:={8000000x_0\over b_0}
 +6(120\cdot10^6x_0)+(120\cdot10^6x_0)^2<2\cdot10^{-15}.
\end{split}                                                   \tag{18.8}
\] Here \(1+10^{-11}/b_0<4\) and \(q_*+8000000x_0/b_0<9\) justify
respectively the square-root bounds two and three used above. The
polynomial monotonicity proved in Section17.3 makes these bounds valid
on the whole interval, not just at its upper endpoint. Define \[
\delta(\xi)=\frac{\beta_N(\xi)+q_*\beta_C(\xi)}{1-\beta_C(\xi)}.
\] Its denominator is positive and
\(\delta(\xi)<(B_N+q_*B_C)/(1-B_C)<10^{-7}\). For every nonzero original
weight, (17.14) permits division by the actual norm. Subtract the two
quotients with their denominators retained: \[
\begin{split}
{w^t\mathcal Nw\over w^tCw}
 -{\kappa\mathcal J_\xi(w)\over\mathcal H_\xi(w)}
&={w^t\mathcal Nw-\kappa\mathcal J_\xi(w)\over w^tCw}\\
&\quad-{\kappa\mathcal J_\xi(w)\over\mathcal H_\xi(w)}
 {w^tCw-\mathcal H_\xi(w)\over w^tCw}.
\end{split}
\] Since \(w^tCw\ge(1-\beta_C)\mathcal H_\xi(w)\), this proves \[
\boxed{\quad
\left|{w^t\mathcal Nw\over w^tCw}
       -{\kappa\mathcal J_\xi(w)\over\mathcal H_\xi(w)}\right|
 \le\kappa\delta(\xi),\quad
0<\xi\le10^{-24},\quad w\ne0.\quad}                         \tag{18.9}
\] The bound is uniform in every original box and every such weight. It
controls all quotients, not only the minimum or the incidence kernel. In
particular they all lie between \(\kappa(3-\delta(\xi))\) and
\(\kappa(q_*+\delta(\xi))\). The independent stronger physical lower
bound \(3\kappa(1-512\xi/3)\) from (15.13) still applies.

For a precise record of the interpolation, put \[
\eta_\xi(w)={b_0\xi^4\|c\|_2^2\over\mathcal H_\xi(w)}.
\] Then \(0\le\eta_\xi(w)\le1\), and the coefficient quotient satisfies
the exact identity \[
\frac{\mathcal J_\xi(w)}{\mathcal H_\xi(w)}
 =3+\frac{87}{49}\eta_\xi(w).                                \tag{18.10}
\] This records both contributions with their original powers of
coupling. It is not a change of state or its norm.

\subsection{18.3. Exact compression, inverse and retained
leakage}\label{exact-compression-inverse-and-retained-leakage}

Let \(\mathcal H_{\rm phys,\mathbb R}\) be the real-valued part of the
gauge-invariant Hilbert space. The real operator \(A\) preserves it. Set
\(\mathscr S_\xi=\operatorname{Ran}V_\xi\subset
 \mathcal H_{\rm phys,\mathbb R}\cap\psi^\perp\), where \(V_\xi w=v_w\).
On the interval above, (17.14) proves \(\dim\mathscr S_\xi=N\). Every
vector of this space is smooth. The adjoint and the orthogonal
projection are \[
\begin{split}
V_\xi^*f&=(\langle v_e,f\rangle)_e,\qquad V_\xi^*V_\xi=C,\\
P_{\mathscr S_\xi}&=V_\xi C^{-1}V_\xi^*.
\end{split}                                                   \tag{18.11}
\] The last operator is self-adjoint because \(C\) is positive and
symmetric. Its square is itself by \(V_\xi^*V_\xi=C\). Its range is
contained in \(\mathscr S_\xi\), and it fixes \(V_\xi w\) for every
\(w\); thus it is the stated projection. The inverse of \(V_\xi\) on its
image is \(C^{-1}V_\xi^*\), which is the exact formula (17.15).

Define the self-adjoint finite-dimensional compression
\(R_\xi=P_{\mathscr S_\xi}A|_{\mathscr S_\xi}\). For
\(u,v\in\mathscr S_\xi\), \(\langle u,R_\xi v\rangle=\langle u,Av\rangle
=\langle Au,v\rangle\). Its coordinate intertwiner and its full
uncompressed action are exactly \[
\begin{split}
R_\xi V_\xi&=V_\xi C^{-1}\mathcal N,\\
AV_\xi&=V_\xi C^{-1}\mathcal N+\mathcal L_\xi,\qquad
\mathcal L_\xi:=(I-P_{\mathscr S_\xi})AV_\xi,\\
V_\xi^*\mathcal L_\xi&=0.
\end{split}                                                   \tag{18.12}
\] These follow by inserting (18.11) and \(V_\xi^*AV_\xi=\mathcal N\).
The domain of \(\mathcal L_\xi\) is the entire finite edge space; its
codomain is \(\mathscr S_\xi^\perp\cap
\mathcal H_{\rm phys,\mathbb R}\cap\psi^\perp\). Thus
\(C^{-1}\mathcal N\) is self-adjoint in the exact inner product
\(\langle w,z\rangle_C=w^tCz\). This records its metric rather than
replacing it by the Euclidean one. The map \(V_\xi\) is an isometry from
this metric space onto \(\mathscr S_\xi\). Its Rayleigh quotient is
exactly the first quotient in (18.9). No vanishing of
\(\mathcal L_\xi\), or invariance of the state image under \(A\), is
used or asserted.

\subsection{18.4. Ordered Ritz values and the exact top coefficient
eigenspace}\label{ordered-ritz-values-and-the-exact-top-coefficient-eigenspace}

Let \(B_0\) be the diagonal coefficient matrix with entries one on faces
and \(b_0\) on pairs, and let \(B_1\) have entries three on faces and
\(1/26\) on pairs. On the original edge space define \[
\mathsf H_\xi=\mathsf T_\xi^tB_0\mathsf T_\xi,\qquad
\mathsf J_\xi=\mathsf T_\xi^tB_1\mathsf T_\xi,
\qquad K^{\rm coeff}_\xi=\mathsf H_\xi^{-1}\mathsf J_\xi.
                                                               \tag{18.13}
\] Their quadratic forms are respectively \(\mathcal H_\xi\) and
\(\mathcal J_\xi\), and (17.6) proves the inverse exists. The operator
\(K^{\rm coeff}_\xi\) is self-adjoint in the positive metric
\(\mathsf H_\xi\). Let its ordered eigenvalues be
\(\mu_1\le\cdots\le\mu_N\), retaining their dependence on \(L,\xi\), and
let \(\lambda_1\le\cdots\le\lambda_N\) be those of \(R_\xi\), with
physical energy units retained.

Here is the full variational comparison argument. For any self-adjoint
operator \(B\) in a finite-dimensional positive inner product \(g\),
take a \(g\)-orthonormal eigenbasis with ordered eigenvalues \(\nu_i\).
A \(j\)-dimensional subspace intersects the \(g\)-orthogonal complement
of the first \(j-1\) eigenvectors in a nonzero vector: the projection
onto their span has rank at most \(j-1\). Such a vector has Rayleigh
quotient at least \(\nu_j\), by its eigenbasis expansion. The span of
the first \(j\) eigenvectors has maximum quotient exactly \(\nu_j\).
Hence \[
\nu_j=\min_{\substack{W\text{ linear}\\\dim W=j}}
       \ \max_{0\ne w\in W}\frac{g(w,Bw)}{g(w,w)}.
                                                               \tag{18.14}
\] The eigenbasis used to prove this identity does not change the
original physical vectors or either retained Gram matrix. Both problems
in (18.13) and (18.12) use precisely the same collection of linear
subspaces of the raw edge space. Apply (18.9) to every nonzero vector of
every one of those subspaces, then take the maximum and minimum in
(18.14). This proves \[
\boxed{\quad
|\lambda_j/\kappa-\mu_j|\le\delta(\xi)
\quad(1\le j\le N),\qquad 3\le\mu_j\le\frac{234}{49}.
\quad}                                                       \tag{18.15}
\] No volume-dependent matrix perturbation constant enters.

The top eigenspace of the coefficient problem is computable exactly,
without solving a large characteristic polynomial: \[
q_*\mathsf H_\xi-\mathsf J_\xi
 =\frac{87}{49}\frac{\xi^2}{16}\mathsf I^t\mathsf I,
\qquad
\ker(K^{\rm coeff}_\xi-q_*I)=\ker\mathsf I .                 \tag{18.16}
\] The first identity follows by retaining both diagonal blocks in
(18.13); the pair block vanishes because \(q_*b_0=1/26\). Its right side
is positive, with zero quadratic form exactly when \(\mathsf Iw=0\).
Invertibility of \(\mathsf H_\xi\) then proves the stated equality of
kernels.

The exact original-box incidence calculation in the parent,
\emph{Electric covariance and unsigned incidence}, Section3,
equations(3.2)--(3.9), gives \(\dim\ker\mathsf I=3m(m+1)=N-M\),
\(m=2L\). For explicitness, its coordinate proof writes
\(w_i(n)=(-1)^{n_1+n_2+n_3}a_i(n)\) and turns the kernel equations into
\(D_ja_i+D_ia_j=0\). Commuting differences in three distinct directions
gives \(D_kD_ja_i=0\). The complete inverse parameters are the
longitudinal boundary values \(f_i(n_i)\) and anchored potentials
\(\varphi_{ij}(x,y)=\sum_{r=-L}^{x-1}u_{ij}(r,y)\), where
\(u_{ij}(x,y)=a_i(x\mathbf e_i+y\mathbf e_j-L\mathbf e_k)-f_i(x)\). They
obey \(\varphi_{ij}(-L,y)=\varphi_{ij}(x,-L)=0\), and reconstruct \[
a_i(n)=f_i(n_i)+\sum_{j>i}D_i\varphi_{ij}(n_i,n_j)
                    -\sum_{j<i}D_i\varphi_{ji}(n_j,n_i).
\] The vanishing transverse mixed difference proves the additive
decomposition giving these parameters; summing \(D_ju_{ij}+D_iu_{ji}=0\)
proves both reconstructed cross derivatives. The boundary values and
finite sums fix the inverse uniquely. There are \(3m\) longitudinal and
\(3m^2\) anchored potential entries, on the original edge and vertex
domains. Thus (18.16), including multiplicities, gives \[
\begin{split}
\mu_{M+1}=\cdots=\mu_N&=\frac{234}{49},\qquad
\mu_j<\frac{234}{49}\quad(1\le j\le M),\\
\left|\lambda_j-\frac{234}{49}\kappa\right|
 &\le\kappa\delta(\xi)\quad(M+1\le j\le N).
\end{split}                                                   \tag{18.17}
\] No uniform separation of \(\mu_M\) from \(q_*\) is asserted. The last
\(N-M\) Ritz values need not be exactly equal; the comparison preserves
the actual mixed energy and covariance.

\subsection{18.5. Consequence for the true physical excitation
eigenvalues}\label{consequence-for-the-true-physical-excitation-eigenvalues}

The gauge-invariant restriction of \(A\) has compact resolvent because
its gauge projection commutes with the full resolvent; the restriction
is a compression of a compact operator. Remove its simple vacuum. The
remaining physical space is infinite dimensional: characters of every
integer and half-integer spin of a fixed face holonomy are gauge
invariant, and are linearly independent by integration in one face link.
Removing one dimension does not change that fact. Let
\(\epsilon_1\le\epsilon_2\le\cdots\) be the true excitation eigenvalues
there, counted with multiplicity.

The eigenbasis proof of (18.14) applies also to this positive
compact-resolvent operator: the first \(j\) eigenvectors give the upper
variational inequality, and expansion of the nonzero intersection with
their first \(j-1\) orthogonal complement gives the lower one. This uses
the form domain for the nonnegative convergent spectral sum. Restricting
the admissible \(j\)-dimensional spaces to
\(\mathscr S_\xi\subset\operatorname{Dom}A\) therefore gives \[
\begin{split}
3\kappa(1-512\xi/3)&\le\epsilon_j\le\lambda_j
 \le\kappa\bigl(\mu_j+\delta(\xi)\bigr)\quad(1\le j\le N),\\
\#\{j:\epsilon_j\le\kappa(q_*+\delta(\xi))\}&\ge N.
\end{split}                                                   \tag{18.18}
\] The lower bound is the independent physical gap proof (15.13). The
upper inequalities and count are for the original interacting operator,
not the free coefficient operator. They do not transfer the exact
coefficient multiplicity in (18.17) into a physical eigenvalue
multiplicity. Equation (18.12) gives the exact map and retained leakage
underlying that distinction. All energies are measured from the actual
\(\mathcal E\); no level below that spectral minimum has been
introduced.

For the original nonnegative native weights in (1.4), the sharper
targeted bounds of Section14 remain in force. The present result
additionally controls every mixed signed weight and the full ordered
electric-state compression. It leaves \(a>0\),
\(\kappa=2g_{\rm YM}^2/a\), \(\xi=1/(4g_{\rm YM}^4)\), and the geometric
input (1.3) unchanged. No spatial continuum limit or Millennium
conclusion is inferred from this finite-regulator strong-coupling
result.

\subsection{18.6. Precise literature dependency and an attainment
counterexample}\label{precise-literature-dependency-and-an-attainment-counterexample}

The canonical corpus routes the variational method to Mathieu Lewin,
\emph{Spectral Theory and Quantum Mechanics}, Universitext,
Springer2024,
DOI\href{https://doi.org/10.1007/978-3-031-66878-4}{10.1007/978-3-031-66878-4},
Section5.6.1, printed pages201--204. Equations(5.42) and(5.44) are the
min--max value and restriction inequality used here; (18.14) and(18.18)
provide their proofs in the exact present spaces. Two details in that
edition must not be imported as extra mathematical claims.

First, the sentence on page202 asserting that the minimizing subspaces
are exactly eigenvector-generated low-energy spaces, and the concluding
necessity argument on page204, fail for the following explicit example.
On \(\mathbb R^3\) with its Euclidean inner product take \[
B=\operatorname{diag}(0,1,2),\qquad
W=\operatorname{span}(e_2,e_1+e_3).
\] For all real \(s,t\), \[
\langle se_2+t(e_1+e_3),B[se_2+t(e_1+e_3)]\rangle
 =s^2+2t^2=\|se_2+t(e_1+e_3)\|^2.                           \tag{18.19}
\] Every nonzero vector of \(W\) has quotient one, the second eigenvalue
of \(B\). The dimension-intersection proof of (18.14) proves this is the
minimizing value. Yet \(W\) is not contained in
\(\operatorname{span}(e_1,e_2)\), and
\(W\cap\operatorname{span}(e_3)=\{0\}\). Failure of containment does not
imply the nonzero intersection used in the source's necessity argument.
The min--max value itself is unaffected; only that characterization of
all attaining spaces is refuted.

Second, the matrix following (5.44) represents the stated ordinary
eigenvalues with an orthonormal basis. For a general basis \(v_i\), the
exact coordinate problem is \(\mathcal Nw=\lambda Cw\), with both Gram
matrices retained, as proved in (18.11)--(18.12). Already on
\(\mathbb R\), the identity operator and basis vector \(v_1=2\) give
\(\mathcal N=(4)\), \(C=(4)\), and eigenvalue one, not four. These
independent finite-dimensional checks explain exactly which source
assertions are used and which are corrected. They assert nothing about
unread editions or other parts of the book. No protected bulk text is
reproduced in this note.

\end{document}